%%%%%%%%%%%%%%%%%%%%%%%%%%%%%%%%%%%%%%%%%
% kaobook
% LaTeX Template
% Version 1.3 (December 9, 2021)
%
% This template originates from:
% https://www.LaTeXTemplates.com
%
% For the latest template development version and to make contributions:
% https://github.com/fmarotta/kaobook
%
% Authors:
% Federico Marotta (federicomarotta@mail.com)
% Based on the doctoral thesis of Ken Arroyo Ohori (https://3d.bk.tudelft.nl/ken/en)
% and on the Tufte-LaTeX class.
% Modified for LaTeX Templates by Vel (vel@latextemplates.com)
%
% License:
% CC0 1.0 Universal (see included MANIFEST.md file)
%
%%%%%%%%%%%%%%%%%%%%%%%%%%%%%%%%%%%%%%%%%

%----------------------------------------------------------------------------------------
%	PACKAGES AND OTHER DOCUMENT CONFIGURATIONS
%----------------------------------------------------------------------------------------

\documentclass[
	a4paper, % Page size
	fontsize=10pt, % Base font size
	twoside=true, % Use different layouts for even and odd pages (in particular, if twoside=true, the margin column will be always on the outside)
	%open=any, % If twoside=true, uncomment this to force new chapters to start on any page, not only on right (odd) pages
	%chapterentrydots=true, % Uncomment to output dots from the chapter name to the page number in the table of contents
	numbers=noenddot, % Comment to output dots after chapter numbers; the most common values for this option are: enddot, noenddot and auto (see the KOMAScript documentation for an in-depth explanation)
]{kaobook}

% Choose the language
\ifxetexorluatex
	\usepackage{polyglossia}
	\setmainlanguage[variant=brazilian]{portuguese}
\else
	\usepackage[brazilian]{babel} % Load characters and hyphenation
\fi
\usepackage{csquotes}	% Quotes in the document language

% Load packages for testing
%\usepackage{showframe} % Uncomment to show boxes around the text area, margin, header and footer
%\usepackage{showlabels} % Uncomment to output the content of \label commands to the document where they are used

% Load the bibliography package
\usepackage{kaobiblio}
\addbibresource{main.bib} % Bibliography file

% Load mathematical packages for theorems and related environments
\usepackage[framed=true]{kaotheorems}

% Load the package for hyperreferences
\usepackage{kaorefs}

% Portuguese strings for the kaorefs reference commands (kaorefs has no Portuguese set)
\renewcommand{\chapternameshort}{Cap.}
\renewcommand{\sectionname}{Seção}
\renewcommand{\sectionnameshort}{Seç.}
\renewcommand{\subsectionname}{Subseção}
\renewcommand{\figurenameshort}{Fig.}
\renewcommand{\tablenameshort}{Tab.}

\graphicspath{{examples/documentation/images/}{images/}} % Paths in which to look for images

% T-SQL listings: kaolst style, SQL Server keywords and accented characters in comments
\lstdefinelanguage{TSQL}[]{SQL}{
	morekeywords={go,exec,execute,procedure,proc,trigger,view,print,declare,
		begin,end,transaction,tran,commit,rollback,return,output,grant,revoke,
		deny,inserted,deleted,top,identity,nvarchar,varchar,char,numeric,money,
		tinyint,bit,datetime,use,database,index,clustered,nonclustered,unique,
		instead,after,if,else,while,exists,raiserror,throw,try,catch,login,
		user,role,schema,with,check,option,cascade,merge,output,over,partition},
	sensitive=false,
}
\lstset{
	style=kaolst,
	language=TSQL,
	inputencoding=utf8,
	extendedchars=true,
	literate={á}{{\'a}}1 {à}{{\`a}}1 {ã}{{\~a}}1 {â}{{\^a}}1 {é}{{\'e}}1
		{ê}{{\^e}}1 {í}{{\'i}}1 {ó}{{\'o}}1 {õ}{{\~o}}1 {ô}{{\^o}}1 {ú}{{\'u}}1
		{ç}{{\c c}}1 {Á}{{\'A}}1 {É}{{\'E}}1 {Ç}{{\c C}}1 {Ã}{{\~A}}1,
}

\makeindex[columns=3, title=Índice Remissivo, intoc] % Make LaTeX produce the files required to compile the index

\makeglossaries % Make LaTeX produce the files required to compile the glossary
\newglossaryentry{computer}{
	name=computer,
	description={is a programmable machine that receives input, stores and manipulates data, and provides output in a useful format}
}

% Glossary entries (used in text with e.g. \acrfull{fpsLabel} or \acrshort{fpsLabel})
\newacronym[longplural={Frames per Second}]{fpsLabel}{FPS}{Frame per Second}
\newacronym[longplural={Tables of Contents}]{tocLabel}{TOC}{Table of Contents}

 % Include the glossary definitions

\makenomenclature % Make LaTeX produce the files required to compile the nomenclature

% Reset sidenote counter at chapters
%\counterwithin*{sidenote}{chapter}

%----------------------------------------------------------------------------------------

\begin{document}

%----------------------------------------------------------------------------------------
%	BOOK INFORMATION
%----------------------------------------------------------------------------------------

\titlehead{%
	\includegraphics[height=2cm]{logo-unicamp.jpg}\hfill
	\includegraphics[height=1.2cm]{logo-extecamp.png}\par\medskip
	Faculdade de Tecnologia --- Unicamp}
\subject{Especialização em Engenharia e Administração de Sistemas de Banco de Dados}

\title{Anotações de Estudo}
\subtitle{Resumos das disciplinas com complementos da bibliografia}

\author{Gabriele Alves}

\date{\today}

%----------------------------------------------------------------------------------------

\frontmatter % Denotes the start of the pre-document content, uses roman numerals

%----------------------------------------------------------------------------------------
%	OPENING PAGE
%----------------------------------------------------------------------------------------

%\makeatletter
%\extratitle{
%	% In the title page, the title is vspaced by 9.5\baselineskip
%	\vspace*{9\baselineskip}
%	\vspace*{\parskip}
%	\begin{center}
%		% In the title page, \huge is set after the komafont for title
%		\usekomafont{title}\huge\@title
%	\end{center}
%}
%\makeatother

%----------------------------------------------------------------------------------------
%	COPYRIGHT PAGE
%----------------------------------------------------------------------------------------

\makeatletter
\uppertitleback{Faculdade de Tecnologia --- Unicamp} % Header

\lowertitleback{
	\textbf{Sobre estas anotações}\\
	Anotações pessoais de estudo. Os resumos foram escritos com palavras
	próprias a partir das aulas, dos slides e da bibliografia de cada
	disciplina, sempre com a fonte indicada. Os slides, scripts e listas de
	exercícios pertencem aos respectivos professores.

	\medskip

	\textbf{Colofão} \\
	Composto em \href{https://www.latex-project.org/}{\LaTeX} com o
	\href{https://sourceforge.net/projects/koma-script/}{\KOMAScript} e a
	classe \href{https://github.com/fmarotta/kaobook/}{kaobook}.

	Código-fonte: \url{https://github.com/gab-i-alves/unicamp-learning-notes}
}
\makeatother

%----------------------------------------------------------------------------------------
%	DEDICATION
%----------------------------------------------------------------------------------------

%----------------------------------------------------------------------------------------
%	OUTPUT TITLE PAGE AND PREVIOUS
%----------------------------------------------------------------------------------------

% Note that \maketitle outputs the pages before here

\maketitle

%----------------------------------------------------------------------------------------
%	PREFACE
%----------------------------------------------------------------------------------------

\chapter*{Prefácio}
\addcontentsline{toc}{chapter}{Prefácio}

Comecei a especialização em Engenharia e Administração de Sistemas de
Banco de Dados da Faculdade de Tecnologia da Unicamp em outubro de 2026.
Estas anotações nasceram da vontade de ter, num lugar só, o que cada aula
ensina e o que a bibliografia e a documentação oficial acrescentam. Os
slides mostram o caminho da professora ou do professor; os livros e os
manuais explicam o porquê e os casos que a aula não teve tempo de cobrir.

Os resumos foram escritos com palavras próprias. Os exemplos também são
próprios: em vez de repetir os bancos dos slides, os diagramas e os
comandos usam outros cenários, como uma biblioteca, para obrigar o
raciocínio a funcionar fora do exemplo original. Quando um trecho vem da
bibliografia ou da documentação, a fonte aparece na margem, ao lado do
texto.

\section*{Como o texto está organizado}

Cada disciplina é uma parte do volume, e cada capítulo acompanha uma aula
ou um bloco de conteúdo. A abertura do capítulo diz de qual aula ele vem
e quais obras o complementam. Quando o livro ou a documentação corrige
ou detalha um ponto dos slides, o trecho vem numa subseção marcada como
complemento, separada do conteúdo da aula.

Os diagramas foram desenhados para estas anotações. Os comandos SQL da
primeira disciplina seguem o dialeto do SQL Server (T-SQL), usado em
aula. Cada disciplina termina com dois capítulos: as avaliações, com o
enunciado resumido, a resposta e as decisões tomadas, e um roteiro de
estudo com perguntas em três níveis, sem as respostas, para revisar o
conteúdo voltando ao texto e às referências. No fim do volume ficam as
referências, o glossário e o índice remissivo.

\section*{Estado atual}

Este é um volume em construção. Até agora, só a CET 615 (Fundamentos de
Sistemas de Banco de Dados) foi concluída: dos conceitos de banco de
dados e SGBD até gatilhos, controle de acesso, as duas avaliações e o
roteiro de estudo. As outras disciplinas entram como novas partes à medida
que o curso avança.

\section*{As disciplinas do curso}

O curso tem 16 disciplinas de 24 horas, distribuídas em quatro blocos de
formação, e um Trabalho de Conclusão de Curso de 64 horas, num total de
448 horas. A lista segue a
\href{https://wordpress.ft.unicamp.br/espbd/?page_id=12}{página oficial de disciplinas}.

\begin{center}
\small
\begin{tabular}{@{}lp{0.62\textwidth}l@{}}
\toprule
\multicolumn{3}{@{}l}{\textbf{Engenharia de Dados}} \\
\midrule
CET 623 & Artificial Intelligence Data Management & \\
CET 615 & Fundamentos de Sistemas de Banco de Dados & concluída \\
CET 602 & Engenharia de Requisitos de Banco de Dados & \\
CET 603 & Projeto de Bancos de Dados & \\
CET 616 & Arquiteturas de Alto Desempenho para BD & \\
CET 619 & BD Aplicados à Internet das Coisas (IoT) & \\
CET 620 & Bancos de Dados Não-relacionais & \\
\midrule
\multicolumn{3}{@{}l}{\textbf{Administração de Banco de Dados (DBA)}} \\
\midrule
CET 604 & Programação e Otimização em SQL & \\
CET 607 & Segurança de Informações & \\
CET 617 & Políticas de Administração de Bancos de Dados & \\
CET 624 & Data Storage Management & \\
\midrule
\multicolumn{3}{@{}l}{\textbf{Ciência de Dados (Análise de Dados)}} \\
\midrule
CET 622 & Data Lake e Data Warehouse & \\
CET 611 & Mineração de Dados & \\
CET 612 & Visualização de Informações & \\
CET 621 & Aprendizado de Máquina na Análise de Dados & \\
\midrule
\multicolumn{3}{@{}l}{\textbf{Arquitetura e Gestão de Dados}} \\
\midrule
CET 618 & Gestão e Governança de Dados & \\
\bottomrule
\end{tabular}
\end{center}

\bigskip
\begin{flushright}
\textit{Gabriele Alves}\\
novembro de 2026
\end{flushright}

\index{preface}

%----------------------------------------------------------------------------------------
%	TABLE OF CONTENTS & LIST OF FIGURES/TABLES
%----------------------------------------------------------------------------------------

\begingroup % Local scope for the following commands

% Define the style for the TOC, LOF, and LOT
%\setstretch{1} % Uncomment to modify line spacing in the ToC
%\hypersetup{linkcolor=blue} % Uncomment to set the colour of links in the ToC
\setlength{\textheight}{230\hscale} % Manually adjust the height of the ToC pages

% Turn on compatibility mode for the etoc package
\etocstandarddisplaystyle % "toc display" as if etoc was not loaded
\etocstandardlines % "toc lines" as if etoc was not loaded

\tableofcontents % Output the table of contents

\listoffigures % Output the list of figures

% Comment both of the following lines to have the LOF and the LOT on different pages
\let\cleardoublepage\bigskip
\let\clearpage\bigskip

\listoftables % Output the list of tables

\endgroup

%----------------------------------------------------------------------------------------
%	MAIN BODY
%----------------------------------------------------------------------------------------

\mainmatter % Denotes the start of the main document content, resets page numbering and uses arabic numbers
\setchapterstyle{kao} % Choose the default chapter heading style

\pagelayout{wide} % No margins
\addpart{CET 615 --- Fundamentos de Sistemas de Banco de Dados}
\pagelayout{margin} % Restore margins

\setchapterpreamble[u]{\margintoc}
\chapter{Conceitos Fundamentais}
\labch{cet615-conceitos}

Este capítulo resume a primeira aula da disciplina\sidecite{BaiocoCET615}:
o que é um banco de dados, o papel do SGBD, as fases do projeto de um
banco e a linguagem SQL. Os complementos vêm de \textcite{Elmasri2011},
6.ª edição. O plano de ensino indica a 7.ª edição; a numeração de
capítulos muda entre as duas, então as referências aqui valem para a 6.ª.

\section{Banco de dados}

Um \textbf{banco de dados} (BD)\index{banco de dados} é uma coleção de dados
relacionados entre si e com significado. Juntar dados ao acaso não cria um
banco: ele é projetado, construído e carregado para um propósito, atende
um grupo de usuários e de aplicações e representa uma parte do mundo real,
o \textbf{minimundo}\index{minimundo} ou universo de discurso. Quando o
minimundo muda, o banco precisa mudar junto.

\section{Do programa ao SGBD}

A aula apresenta três formas de guardar dados em disco, em ordem de
evolução:

\begin{enumerate}
	\item \textbf{Dados dentro do programa.} Código e dados formam um único
	objeto; nenhum outro programa consegue usar os dados.
	\item \textbf{Sistema de arquivos.} Os dados ficam em arquivos separados
	dos programas, mas cada programa carrega no próprio código as rotinas de
	gravar e ler. Mudar o formato de um arquivo obriga a mudar todos os
	programas que o usam.
	\item \textbf{Sistema de gerenciamento de banco de dados.} O SGBD fica
	entre os programas e os arquivos e cuida do armazenamento e da
	recuperação. Os programas ficam independentes da forma como os dados
	estão gravados.
\end{enumerate}

Um \textbf{SGBD}\index{SGBD} é um software de uso geral para definir,
construir e manipular bancos de dados de aplicações diferentes. As
capacidades básicas citadas em aula são:

\begin{itemize}
	\item compartilhamento de dados com controle de concorrência;
	\item controle de acesso para vários usuários;
	\item garantia das restrições de integridade, como tipos de dados e
	relacionamentos;
	\item cópia de segurança (\emph{backup}) e recuperação
	(\emph{recovery}) depois de falhas de hardware ou software.
\end{itemize}

\subsection*{Complemento: o que distingue a abordagem de banco de dados}

\textcite[seção 1.3]{Elmasri2011} descreve quatro características que
separam o banco de dados do sistema de arquivos:

\begin{description}
	\item[Natureza autodescritiva.] O SGBD guarda, no
	\textbf{catálogo}\index{catálogo}, a definição completa do banco:
	tabelas, tipos, restrições. Esses são os \textbf{metadados}. Um
	programa sobre arquivos depende de definições escritas no próprio
	código.
	\item[Isolamento entre programas e dados.] A estrutura dos dados fica
	no catálogo, e não no programa. Por isso é possível mudar a estrutura
	sem reescrever os programas. O livro chama essa propriedade de
	\emph{program-data independence}.
	\item[Múltiplas visões.] Usuários diferentes podem ver subconjuntos
	diferentes do banco, inclusive dados derivados que não estão gravados.
	As visões do \refch{cet615-visoes} são a implementação disso em SQL.
	\item[Compartilhamento e transações multiusuário.] Vários usuários
	acessam o banco ao mesmo tempo, e o SGBD controla a concorrência para
	que cada transação veja um estado consistente.
\end{description}

\section{Sistema de banco de dados}

O \textbf{sistema de banco de dados}\index{sistema de banco de dados}
(SBD) é o conjunto formado pelos bancos de dados e pelo SGBD. A aula
divide o SBD em três partes:

\begin{description}
	\item[Interface com o usuário:] valida os pedidos dos usuários e
	inicia as operações correspondentes no banco.
	\item[Mecanismo de acesso:] mantém as estruturas de dados, acessa os
	dados e garante a integridade.
	\item[Armazenamento:] os arquivos que contêm os dados.
\end{description}

E cita quatro perfis de profissionais:

\begin{description}
	\item[Administrador do banco (DBA):] autoriza o acesso, coordena e
	monitora o uso e resolve problemas como falhas de segurança e baixo
	desempenho.
	\item[Projetista do banco:] decide quais dados guardar e escolhe as
	estruturas que os representam.
	\item[Usuário final:] consulta e altera dados e gera relatórios.
	\item[Analistas e programadores:] os analistas especificam as
	transações que atendem aos usuários; os programadores as implementam.
\end{description}

\subsection*{Complemento: arquitetura de três esquemas}

A aula não trata deste ponto, mas ele explica de onde vem a
independência de dados. \textcite[seção 2.2]{Elmasri2011} descreve três
níveis de esquema:

\begin{description}
	\item[Nível interno:] como os dados estão gravados fisicamente
	(arquivos, índices, caminhos de acesso).
	\item[Nível conceitual:] a estrutura do banco inteiro para a
	comunidade de usuários: entidades, tipos, relacionamentos e
	restrições, sem detalhes de armazenamento.
	\item[Nível externo:] as visões de cada grupo de usuários.
\end{description}

O SGBD converte pedidos de um nível para o outro por meio de
\emph{mapeamentos}. Daí saem dois tipos de independência:

\begin{itemize}
	\item \textbf{Independência lógica:}\index{independência de dados}
	mudar o esquema conceitual (por exemplo, criar uma coluna) sem mudar
	as visões nem os programas que não usam essa mudança.
	\item \textbf{Independência física:} mudar o esquema interno (por
	exemplo, criar um índice ou reorganizar arquivos) sem mudar o esquema
	conceitual.
\end{itemize}

Segundo o livro, a independência física é comum na prática; a lógica é
mais difícil, porque os programas dependem da estrutura que consultam.

\section{Projeto de banco de dados}

Quatro termos aparecem em todo o projeto:

\begin{description}
	\item[Modelo de dados:] conjunto de conceitos e regras para descrever a
	estrutura de um banco (por exemplo, o modelo entidade-relacionamento ou
	o modelo relacional).
	\item[Modelagem:] o processo de estruturar a informação do minimundo
	segundo um modelo de dados.
	\item[Esquema:] a descrição do banco, em texto ou diagrama, conforme um
	modelo. Muda pouco.
	\item[Instância (ou estado):] os dados do banco em um momento. Muda a
	cada inserção, alteração ou exclusão.
\end{description}

Dois alertas feitos em aula\sidecite{BaiocoCET615Aulas}:

\begin{itemize}
	\item O resultado da modelagem é um \emph{esquema}, não ``o modelo''. O
	modelo é a linguagem comum que permite a todos ler o esquema do mesmo
	jeito. A expressão ``o modelo do meu banco'' é comum no mercado, mas
	troca os conceitos.
	\item O esquema não se deduz das instâncias. Olhando os dados, uma
	relação pode parecer 1:N e não ser. Quando o SSMS desenha o diagrama
	de um banco pronto, ele lê a definição guardada no catálogo, e não os
	dados.
\end{itemize}

\subsection{Fases do projeto}

\begin{enumerate}
	\item \textbf{Coleta e análise de requisitos.} Resulta em requisitos de
	dados e requisitos funcionais.
	\item \textbf{Projeto conceitual.} Modela os requisitos de dados no
	ME-R (ou ME-R estendido) e gera o \emph{esquema conceitual},
	independente de SGBD. Os requisitos funcionais viram a especificação
	das transações.
	\item \textbf{Projeto lógico.} Mapeia o esquema conceitual para o
	modelo do tipo de SGBD escolhido (aqui, o relacional) e gera o
	\emph{esquema lógico}. Depende do tipo de SGBD, mas não do fabricante.
	\item \textbf{Projeto físico.} Implementa o esquema lógico com as
	estruturas de armazenamento e acesso de um SGBD específico, usando a
	DDL dele. Gera o \emph{esquema físico}.
\end{enumerate}

\section{Projeto conceitual: o ME-R}

\marginnote{Em ordem: o modelo relacional de Codd é de 1970; o ME-R de
Peter Chen, de 1976. O ME-R surgiu porque os esquemas relacionais
cresceram e faltava um modelo para conceber o banco antes de
implementá-lo.}
O modelo entidade-relacionamento\index{ME-R} tem quatro construtores:
conjunto de entidades, conjunto de relacionamentos, atributos de
entidades e atributos de relacionamentos.

\begin{description}
	\item[Entidade:] algo sobre o qual se quer guardar informação. Pode ser
	físico (pessoa, carro) ou conceitual (departamento, curso).
	\item[Conjunto de entidades (CE):] agrupa entidades com propriedades
	semelhantes. No diagrama, um retângulo.
	\item[Relacionamento:] associação entre entidades, como
	\emph{Funcionário trabalha em Departamento}.
	\item[Conjunto de relacionamentos (CR):] agrupa relacionamentos
	semelhantes. No diagrama, um losango.
\end{description}

\subsection{Grau e cardinalidade}

O \textbf{grau} é o número de conjuntos de entidades que participam do
relacionamento: binário (dois), ternário (três) e assim por diante. Um
relacionamento pode ligar um conjunto a ele mesmo, como \emph{Disciplina
é pré-requisito de Disciplina}; nesse caso, os papéis (``é pré-requisito
de'' e ``pré-requisita'') distinguem as duas pontas.

A \textbf{cardinalidade}\index{cardinalidade} diz com quantas entidades
do outro lado cada entidade pode se relacionar:

\begin{itemize}
	\item \textbf{1:1}: um funcionário gerencia no máximo um departamento,
	e cada departamento tem no máximo um gerente.
	\item \textbf{1:N}: um funcionário pode ter vários dependentes, e cada
	dependente depende de um só funcionário.
	\item \textbf{N:M}: um funcionário desenvolve vários projetos, e cada
	projeto tem vários funcionários.
\end{itemize}

Nos ternários, a cardinalidade combina as três pontas: 1:1:1, 1:1:N,
1:N:M ou N:M:P. Para ler um ternário, fixam-se duas pontas e pergunta-se
pela terceira: dado um aluno e um professor supervisor, quantas
disciplinas esse aluno monitora? A resposta vem dos requisitos.

Em aula, a professora completou o slide com a \textbf{multiplicidade} de
cada papel, que tem um mínimo (0 ou 1) e um máximo (1 ou N). A
cardinalidade do relacionamento é formada pelos \emph{máximos}, e é ela
que define o mapeamento para o relacional. O mínimo só diz se a
participação é obrigatória: pode existir departamento sem funcionário?
Também não há sentido de leitura no diagrama: 1:N e N:1 são o mesmo
relacionamento visto de lados opostos, e usar a mesma letra nos dois
lados (N e N) não significa quantidades iguais.

\subsection{Atributos e chaves}

Atributos descrevem entidades e relacionamentos; seus valores são os
dados guardados no banco. Na convenção da disciplina, atributos de
entidades são obrigatórios e atributos de relacionamentos são opcionais.

A \textbf{chave}\index{chave} de um conjunto de entidades é o conjunto
de atributos cujos valores identificam cada entidade de forma única. Ela
pode ser composta: uma turma pode ser identificada por sigla e período
juntos. O ME-R indica uma só chave por conjunto de entidades, e
relacionamentos não têm chave. Um atributo de relacionamento é o que só
faz sentido na associação: a nota pertence ao par aluno--disciplina, não
ao aluno nem à disciplina.

\section{ME-R estendido}

Modelar é escolher entre abstrair (agrupar em algo mais genérico) e
detalhar. O ME-R estendido acrescenta duas abstrações ao ME-R:
generalização e agregação.

\subsection{Generalização}

A \textbf{generalização}\index{generalização} reúne as características
comuns de vários conjuntos de entidades em um conjunto genérico. No
sentido contrário, a especialização detalha o genérico em conjuntos
específicos. A relação entre eles é ``é um'', e os específicos herdam os
atributos do genérico.

Segundo a aula, vale usar generalização em dois casos:

\begin{enumerate}
	\item quando alguns atributos só se aplicam a parte das entidades do
	genérico (só o professor tem área de pesquisa);
	\item quando um relacionamento só envolve parte delas (só o professor
	ministra disciplina).
\end{enumerate}

Basta um dos dois motivos para justificar um específico. Um específico
também pode ter chave própria (o ISBN de um livro, numa livraria); no
relacional, ela vira chave secundária, porque a primária é sempre a
herdada do genérico.

Duas propriedades classificam uma generalização, e elas são
independentes:

\begin{description}
	\item[Participação total ou parcial:] na total, toda entidade do
	genérico pertence a pelo menos um específico; na parcial, pode não
	pertencer a nenhum.
	\item[Exclusão mútua ou sobreposição:] na exclusão mútua, uma entidade
	pertence a no máximo um específico; na sobreposição, pode pertencer a
	vários.
\end{description}

Na notação da disciplina, a haste dupla entre o genérico e o triângulo
indica participação total, e a simples indica parcial. O triângulo vazio
indica exclusão mútua, e o pintado, sobreposição.

\marginnote{Correspondência de termos: participação total ou parcial =
\emph{completeness constraint}; exclusão mútua = \emph{disjoint}
(\texttt{d} no diagrama); sobreposição = \emph{overlapping} (\texttt{o}).}
\textcite[seção 8.3]{Elmasri2011} usa os mesmos dois eixos com outros
nomes e chega às mesmas quatro combinações. O livro observa ainda que,
quando a especialização é definida pelo valor de um atributo (como o
\texttt{tipo} das obras literárias no exemplo da avaliação), a exclusão
mútua vem de graça, porque o atributo só tem um valor por entidade.

\subsection{Agregação}

A \textbf{agregação}\index{agregação} transforma um relacionamento, com
as entidades que ele liga, em um objeto agregado que pode ter atributos
próprios e participar de outros relacionamentos. A aula cita dois casos
de uso:

\begin{enumerate}
	\item representar um \emph{relacionamento entre relacionamentos}: o
	trabalho de um empregado em um projeto usa máquinas;
	\item dar uma chave própria a cada ocorrência do relacionamento: cada
	venda de um produto por um vendedor a um cliente recebe um código. A
	alternativa é uma chave adicional, como a data, somada às chaves dos
	participantes.
\end{enumerate}

A professora insistiu que \emph{nem todo relacionamento vira agregação}.
Um funcionário se associa a um dependente uma vez só, e isso não pede
agregação. Ela faz sentido quando a mesma associação se repete em
ocasiões diferentes e todas precisam ser guardadas: sem a agregação, uma
segunda venda do mesmo produto, pelo mesmo vendedor, ao mesmo cliente
apenas sobrescreveria a data da primeira.

Sobre a escolha da chave:

\begin{itemize}
	\item se os requisitos já trazem um identificador, como o número da
	nota fiscal, use-o; senão, gere um código;
	\item a chave adicional (chaves das partes mais a data) basta quando o
	agregado não se relaciona com mais nada;
	\item se o agregado vai se relacionar com outro conjunto, prefira a
	chave própria e mantenha a combinação das partes como restrição de
	unicidade.
\end{itemize}

No desenho, o retângulo da agregação envolve o losango sem se fundir a
ele, para que se veja quais hastes ligam ao relacionamento e quais ligam
ao agregado. A agregação sempre envolve um único relacionamento.

\subsection*{Complemento: agregação no Elmasri}

\textcite[seção 8.7]{Elmasri2011} trata a agregação como um conceito
geral de abstração, e não como construtor do diagrama EER. O livro
discute o mesmo problema do primeiro caso acima: o ER não permite
relacionamento entre relacionamentos. A solução que ele propõe dentro do
ER é criar uma entidade fraca para o relacionamento e ligá-la ao outro
conjunto. O agregado do ME-R estendido da disciplina cumpre esse papel
direto no diagrama. Ao mapear para o modelo relacional, os dois caminhos
geram a mesma tabela: a tabela \texttt{compra} da avaliação é um exemplo.

\section{Projeto lógico: o modelo relacional}

O modelo relacional\index{modelo relacional} tem dois construtores:
\textbf{relações} e \textbf{atributos}. Informalmente, relação é tabela,
atributo é coluna e tupla é linha (ou registro). Um esquema lógico é
escrito como \texttt{Empregado = \{NumEmp, NomeEmp, Salario, Depto\}},
com a chave primária sublinhada.

\subsection{Chaves}

\begin{description}
	\item[Chave:] subconjunto mínimo de atributos cujos valores são únicos
	em cada tupla. Mínimo quer dizer que, se qualquer atributo sair, a
	unicidade se perde.
	\item[Chave primária:] a chave escolhida para identificar as tuplas.
	Na aula, o critério é a chave mais acessada. As outras chaves são as
	\emph{candidatas} (ou secundárias): em \texttt{Pessoa}, se o CPF é a
	primária, o RG é candidata.
	\item[Chave estrangeira:] atributo cujo valor referencia a chave
	primária de outra relação e cria o vínculo entre as duas. No esquema
	lógico, ela é anotada assim: \emph{Depto -- chave estrangeira
	referenciando Departamento}.
\end{description}

\subsection{Restrições de integridade}

São regras sobre os valores que as relações podem guardar:

\begin{description}
	\item[Integridade de chave:] o valor da chave não se repete entre as
	tuplas.
	\item[Integridade de entidade:] a chave primária nunca é nula.
	\item[Integridade referencial:] o valor de uma chave estrangeira existe
	como chave primária na relação referenciada (ou é nulo, se o atributo
	aceita nulo).
\end{description}

A professora destacou um ponto que costuma confundir: a chave primária
nunca pode ser nula, mas a chave estrangeira e a chave secundária podem.
Um gerente que não trabalha em nenhum departamento tem a chave
estrangeira nula, e isso não viola nada.

\textcite[seção 3.2]{Elmasri2011} acrescenta as \textbf{restrições de
domínio} (cada valor pertence ao tipo e ao intervalo do atributo) e
classifica as três da aula como restrições expressas no próprio esquema,
diferentes das regras de negócio, que precisam de código. No SQL Server,
essas regras de negócio aparecem como \emph{check}, gatilhos e
procedimentos (\refch{cet615-gatilhos}).

\subsection{Mapeamento do ME-R estendido para o relacional}

A aula mostra o mapeamento com o sistema de venda de obras literárias
usadas, o mesmo esquema da primeira avaliação:

\begin{itemize}
	\item cada conjunto de entidades vira uma relação (\texttt{cliente},
	\texttt{obras\_literarias});
	\item o relacionamento 1:N \emph{Vender} não vira relação: a chave do
	lado 1 entra como chave estrangeira do lado N
	(\texttt{login\_vendedor} em \texttt{obras\_literarias}), e o atributo
	do relacionamento (\texttt{data}) vai junto;
	\item a generalização vira uma relação para o genérico, com um
	atributo \texttt{tipo}, e uma relação por específico, cuja chave
	primária é também chave estrangeira para o genérico (\texttt{livro},
	\texttt{periodico});
	\item a agregação \emph{Compra} vira uma relação com chave própria
	(\texttt{codcompra}) e chaves estrangeiras para o comprador e para a
	obra.
\end{itemize}

Dois detalhes do mapeamento só aparecem no projeto físico:

\begin{itemize}
	\item chaves secundárias (como o RG do cliente) não saem
	automaticamente do mapeamento; é preciso revisá-las e declará-las;
	\item dentro da agregação, a obra participa de no máximo uma compra (o
	lado 1). Uma chave estrangeira comum aceita repetição, então a
	garantia vem de um índice \emph{único} sobre ela. É o papel do
	\texttt{create unique index ix\_compra\_obra on compra(codobra)} no
	esquema da avaliação.
\end{itemize}

No modelo relacional, a chave estrangeira referencia a chave primária
da outra relação. O SQL Server também aceita referenciar uma coluna
\texttt{unique}; a professora chama isso de ``burlar'' o modelo.

Em aula, a professora justificou o atributo \texttt{tipo}: com exclusão
mútua, ele permite saber o específico de cada obra consultando só a
tabela genérica. Com sobreposição, um único \texttt{tipo} não basta.

\subsection*{Complemento: as quatro opções de mapeamento}

\textcite[seção 9.2, passo 8]{Elmasri2011} lista quatro opções para
mapear especializações:

\begin{description}
	\item[8A -- genérico e específicos em relações separadas.] Funciona
	para qualquer combinação de participação e exclusão. É a opção da
	aula.
	\item[8B -- só os específicos.] Cada específico recebe também os
	atributos do genérico. Exige participação total e só é recomendada com
	exclusão mútua, senão a mesma entidade se repete em várias relações.
	\item[8C -- uma relação com um atributo de tipo.] Todos os atributos
	ficam numa tabela só. Exige exclusão mútua e pode gerar muitos nulos.
	\item[8D -- uma relação com um atributo booleano por específico.]
	Serve para sobreposição.
\end{description}

O esquema da avaliação combina a 8A com o atributo de tipo da 8C: as
tabelas ficam separadas, e o \texttt{tipo} evita juntar com cada
específico só para saber de qual se trata.

\section{A linguagem SQL}

A SQL\index{SQL} é a linguagem padrão dos SGBDs relacionais. A ANSI
publicou o primeiro padrão em 1986 (SQL-86), adotado pela ISO em 1987.
Vieram depois a SQL-89 e a SQL-92; a SQL:1999 (SQL3), com recursos
objeto-relacionais; a SQL:2003, com XML; a SQL:2006, com mais integração
a XML; e as revisões de 2008, 2011 e 2016.

A aula divide a SQL em três sublinguagens:

\begin{description}
	\item[DDL (definição):] cria, altera e remove a estrutura e os
	objetos do banco: tabelas, índices, visões, procedimentos e gatilhos
	(\refch{cet615-ddl}).
	\item[DML (manipulação):] insere, consulta, altera e exclui dados
	(\refch{cet615-dml}).
	\item[DCL (controle):] usuários, permissões, segurança e
	administração (\refch{cet615-dcl}).
\end{description}

\textcite[cap. 4]{Elmasri2011} explica que, a partir da SQL:1999, o
padrão se divide em um núcleo, que todo SGBD compatível deveria
implementar, e em extensões opcionais para áreas específicas. Os
fabricantes somam a isso recursos próprios. O T-SQL (Transact-SQL), usado
no resto desta parte, é o dialeto do SQL Server.

\setchapterpreamble[u]{\margintoc}
\chapter{O SQL Server e o Banco de Dados}
\labch{cet615-sqlserver}

Este capítulo cobre a segunda apresentação da primeira
aula~\cite{BaiocoCET615}: a ferramenta de administração, os bancos que o
próprio SQL Server mantém, os objetos de um banco e os arquivos em disco.
Os complementos vêm da documentação oficial no Microsoft Learn.

\section{SQL Server Management Studio}

O \textbf{SSMS}\index{SSMS} é a interface gráfica para administrar o SQL
Server desde a versão 2005; antes dela, o papel era do \emph{Enterprise
Manager}. Quase tudo o que se faz em T-SQL também pode ser feito pelo SSMS,
e o editor de consultas do SSMS executa os comandos T-SQL diretamente.

\section{Bancos do sistema e bancos de usuário}

Um servidor SQL Server guarda dois tipos de banco: os \textbf{bancos do
sistema}, que o SGBD usa para se controlar, e os \textbf{bancos de
usuário}, criados para as aplicações. Os bancos do sistema guardam
metadados, os ``dados sobre os dados''.

\begin{description}
	\item[\texttt{master}:]\index{master@\texttt{master}} registra o nível de
	sistema da instância: logins, bancos existentes, espaço de cada um,
	configurações. A aula recomenda manter sempre um backup atualizado dele.
	\item[\texttt{model}:] serve de molde para os bancos novos. Tudo o que
	estiver no \texttt{model} (opções, tamanhos, permissões padrão, tipos)
	é copiado para cada banco criado depois.
	\item[\texttt{tempdb}:] área de trabalho para tabelas temporárias e
	resultados intermediários de consultas, como as que usam
	\texttt{group by} e \texttt{order by}.
	\item[\texttt{msdb}:] usado pelo SQL Server Agent para agendar tarefas
	(\emph{jobs}) e alertas.
\end{description}

Em aula\sidecite{BaiocoCET615Aulas}, a professora reforçou que o SQL
Server \emph{distribui} o catálogo: o que vale para a instância toda
(logins, lista de bancos) fica no \texttt{master}, e o esquema de cada
banco fica nas tabelas de sistema do próprio banco. Sem o \texttt{master},
os arquivos no disco não podem ser interpretados. Essa divisão é do SQL
Server; outros SGBDs, como o Oracle, organizam o catálogo de outro jeito.

Ela mostrou isso ao vivo: \texttt{select * from sys.databases}, executado
no \texttt{master}, devolve a mesma lista de bancos que o SSMS exibe, com
mais colunas. Todo comando DDL apenas grava no catálogo; depois de um
\texttt{create database}, aparece uma linha nova em \texttt{sys.databases}.
O SSMS é ``uma casca'' sobre esse catálogo, e o banco continua existindo
sem a ferramenta. Um efeito prático: o SSMS só relê o catálogo quando se
pede, então um banco criado por script só aparece no \emph{Object
Explorer} depois de um \emph{Refresh}.

\subsection*{Complemento: o banco \texttt{Resource} e a vida do \texttt{tempdb}}

A documentação lista um quinto banco do sistema que não aparece no SSMS:
o \texttt{Resource}\sidecite{MSSystemDatabases}. Ele é somente leitura e
guarda fisicamente os objetos de sistema que acompanham o produto; esses
objetos aparecem, logicamente, no esquema \texttt{sys} de cada banco.

Sobre o \texttt{tempdb}, a documentação\sidecite{MSTempdb} detalha dois
momentos de limpeza diferentes:

\begin{itemize}
	\item tabelas temporárias locais (\texttt{\#tabela}) e procedimentos
	temporários somem quando a sessão que os criou se desconecta;
	\item o \texttt{tempdb} inteiro é recriado vazio toda vez que o serviço
	do SQL Server inicia. Por isso não se faz backup dele.
\end{itemize}

Além das tabelas temporárias, o \texttt{tempdb} guarda objetos internos,
como as tabelas de trabalho de ordenações e agrupamentos, e o
\emph{version store}, usado pelo isolamento por versão de linha e pelos
gatilhos \texttt{after} (\refch{cet615-gatilhos}).

\section{Tabelas do sistema e procedimentos do sistema}

Os slides listam tabelas de sistema presentes em cada banco
(\texttt{sysobjects}, \texttt{syscolumns}, \texttt{sysindexes},
\texttt{sysforeignkeys}, \texttt{systypes}, \texttt{sysusers}) e outras só
no \texttt{master} (\texttt{sysdatabases}, \texttt{syslogins},
\texttt{sysmessages}, \texttt{sysprocesses}, \texttt{sysconfigures}).

Os \textbf{procedimentos do sistema}\index{procedimento do sistema}, com
prefixo \texttt{sp\_}, automatizam tarefas comuns:

\begin{lstlisting}
exec sp_help obras_literarias   -- estrutura de um objeto
exec sp_helpdb espbd            -- dados de um banco e seus arquivos
exec sp_helpindex compra        -- índices de uma tabela
exec sp_helptext ins_livro      -- código de visão, procedimento ou gatilho
\end{lstlisting}

\subsection*{Complemento: visões de compatibilidade e visões de catálogo}

Desde o SQL Server 2005, as tabelas \texttt{sys...} do SQL Server 2000
existem só como \textbf{visões de compatibilidade}\sidecite{MSCompatViews}.
Elas mostram apenas os metadados que existiam na versão 2000 e podem
devolver nulo ou estourar em colunas de identificadores. Para recursos
novos, a documentação recomenda as \textbf{visões de catálogo}\index{visão
de catálogo}:

\begin{table}[h]
	\begin{tabular}{ll}
		\toprule
		Compatibilidade & Catálogo \\
		\midrule
		\texttt{sysobjects} & \texttt{sys.objects} \\
		\texttt{syscolumns} & \texttt{sys.columns} \\
		\texttt{sysindexes} & \texttt{sys.indexes} \\
		\texttt{sysforeignkeys} & \texttt{sys.foreign\_keys} \\
		\texttt{systypes} & \texttt{sys.types} \\
		\texttt{sysusers} & \texttt{sys.database\_principals} \\
		\texttt{sysdatabases} & \texttt{sys.databases} \\
		\texttt{sysprocesses} & \texttt{sys.dm\_exec\_requests} \\
		\bottomrule
	\end{tabular}
	\caption{Visões de compatibilidade e as visões de catálogo que as substituem.}
\end{table}

A documentação também desaconselha alterar objetos de sistema
diretamente: a administração deve passar pelo SSMS, pela DDL e pelos
procedimentos do sistema\sidecite{MSSystemDatabases}.

\section{Objetos de um banco}

Cada banco contém objetos (componentes lógicos):

\begin{description}
	\item[Tabelas:] as relações do modelo relacional.
	\item[Índices:] aceleram o acesso às linhas das tabelas.
	\item[Visões:] uma forma alternativa de ver dados de uma ou mais
	tabelas (\refch{cet615-visoes}).
	\item[Procedimentos armazenados:] conjuntos de comandos SQL compilados
	e guardados no servidor (\refch{cet615-procedimentos}).
	\item[Gatilhos:] procedimentos especiais que o SGBD executa sozinho
	quando há \texttt{insert}, \texttt{update} ou \texttt{delete} na tabela
	a que estão ligados (\refch{cet615-gatilhos}).
	\item[Outros:] tipos de dados, restrições, \emph{defaults} e regras.
\end{description}

\section{T-SQL}

O \textbf{Transact-SQL}\index{T-SQL} é a implementação da SQL no SQL
Server: segue o padrão ANSI/ISO e acrescenta extensões, como variáveis,
controle de fluxo e procedimentos. A divisão da aula para o T-SQL é:

\begin{description}
	\item[DDL:] criar, alterar e remover bancos, tabelas, índices, visões,
	procedimentos, gatilhos, tipos, restrições, \emph{defaults} e regras.
	\item[DML:] \texttt{select}, \texttt{insert}, \texttt{update},
	\texttt{delete}, transações e a execução de visões e procedimentos.
	\item[DCL:] logins, usuários e grupos, privilégios, backup e
	recuperação, \emph{jobs}.
\end{description}

Cada fabricante tem sua extensão procedural: T-SQL no SQL Server, PL/SQL
no Oracle, SQL PL no DB2. A SQL é declarativa (diz \emph{o que} se quer);
os comandos procedurais são imperativos (dizem \emph{como} fazer). Mesmo
comandos padronizados variam entre fabricantes, por exemplo nos tipos de
dados. O \texttt{create trigger} existe no padrão, mas o corpo do gatilho
é escrito na extensão de cada um.

\subsection*{Complemento: lotes e o comando \texttt{go}}

\index{go@\texttt{go}}Os scripts da disciplina separam trechos com
\texttt{go}. Em aula, a professora disse que ele não é necessário no SSMS
e que existe para utilitários de linha de comando. A documentação
precisa esse ponto\sidecite{MSGo}: \texttt{go} não é um comando T-SQL e
nunca chega ao servidor, mas é reconhecido tanto pelo \texttt{sqlcmd}
quanto pelo editor do SSMS. Ele marca o fim de um \textbf{lote}
(\emph{batch}): tudo o que vem antes é enviado e compilado junto.

Por isso o \texttt{go} faz diferença mesmo no SSMS quando se executa um
script inteiro:

\begin{itemize}
	\item variáveis locais só existem dentro do lote em que foram
	declaradas;
	\item \texttt{create view}, \texttt{create procedure} e
	\texttt{create trigger} precisam ser o primeiro comando do lote;
	\item fora da primeira linha do lote, a chamada de um procedimento
	precisa de \texttt{exec}.
\end{itemize}

Ao executar só um trecho selecionado, como a professora faz em aula, o
trecho já é um lote, e o \texttt{go} não muda nada.

\section{Arquivos de um banco}

Cada banco corresponde a arquivos no disco:

\begin{description}
	\item[Arquivo de dados primário (\texttt{.mdf}):] contém as informações
	de inicialização do banco, aponta para os demais arquivos e guarda as
	tabelas do sistema. Todo banco tem exatamente um.
	\item[Arquivos de dados secundários (\texttt{.ndf}):] opcionais;
	acrescentados conforme o banco cresce.
	\item[Arquivo de log de transações (\texttt{.ldf}):] registra as
	operações feitas no banco. Pode haver mais de um. A regra prática da
	aula é dimensionar o log entre 10\% e 15\% do tamanho do banco (em
	aula, ela falou em até 20\%).
\end{description}

O log não acompanha o tamanho do banco porque cresce e encolhe conforme
as transações terminam e o backup de log é feito. Um log que cresce
demais indica transações longas ou backup de log que não está sendo
feito. A professora também chamou a atenção para a separação entre a
visão lógica e a física: o usuário vê tabelas, mas no disco há poucos
arquivos, não um por tabela. Pelo arquivo não se descobre a estrutura do
banco nem se ele está cheio, porque o tamanho é reservado antes. Isso
soma à segurança, junto com a criptografia.

\subsection*{Complemento: grupos de arquivos}

A documentação\sidecite{MSFilesFilegroups} acrescenta pontos que a aula
não trata:

\begin{itemize}
	\item as extensões \texttt{.mdf}, \texttt{.ndf} e \texttt{.ldf} são
	recomendadas, não obrigatórias;
	\item arquivos de dados se organizam em \textbf{grupos de arquivos}
	(\emph{filegroups})\index{filegroup}. O grupo \texttt{PRIMARY} contém o
	arquivo primário e é o grupo padrão, a menos que outro seja definido
	com \texttt{alter database}. Uma tabela pode ser criada em um grupo
	específico (\texttt{create table ... on grupo});
	\item arquivos de log nunca pertencem a grupo de arquivos;
	\item dentro de um grupo, o SQL Server distribui as gravações entre os
	arquivos em proporção ao espaço livre de cada um;
	\item a recomendação é separar dados e log em discos diferentes; a
	maioria dos bancos funciona bem com um arquivo de dados e um de log.
\end{itemize}

Cada arquivo tem dois nomes: o \textbf{lógico} (\texttt{name}), usado nos
comandos T-SQL, e o \textbf{físico} (\texttt{filename}), o caminho no
sistema operacional.

\section{Criar, alterar e excluir bancos}

Os três comandos são executados a partir do banco \texttt{master}.

\subsection{\texttt{create database}}

\begin{lstlisting}[caption={Banco com um arquivo de dados e um de log.}]
use master
go

create database biblioteca
on
   (name = biblioteca_dados,
    filename = 'C:\dados\biblioteca.mdf',
    size = 10MB,
    maxsize = 50MB,
    filegrowth = 5MB)
log on
   (name = biblioteca_log,
    filename = 'C:\dados\biblioteca.ldf',
    size = 2MB,
    maxsize = 10MB,
    filegrowth = 10%)
\end{lstlisting}

Observações da aula sobre o comando:

\begin{itemize}
	\item \texttt{size} e \texttt{maxsize} aceitam KB, MB, GB e TB; sem
	unidade, vale MB;
	\item \texttt{filegrowth} pode ser absoluto (\texttt{5MB}) ou
	percentual (\texttt{10\%});
	\item não se usam decimais: em vez de \texttt{2.5MB}, escreve-se
	\texttt{2560KB};
	\item uma instância aceita no máximo 32.767 bancos;
	\item o caminho em \texttt{filename} depende da instalação; o script
	da aula usa a pasta \texttt{DATA} da instância.
\end{itemize}

\subsection*{Complemento: o que é obrigatório no \texttt{create database}}

Em aula, a professora apresentou o \texttt{log on} como obrigatório e o
arquivo secundário como opcional (normalmente acrescentado depois, com
\texttt{alter database}). A documentação\sidecite{MSCreateDatabase} diz
que todo banco tem pelo menos dois arquivos, um de dados primário e um
de log, mas a \emph{cláusula} \texttt{log on} é opcional: sem ela, o SQL
Server cria sozinho um arquivo de log com 25\% da soma dos arquivos de
dados ou 512 KB, o que for maior. Sem nenhuma especificação de arquivo,
o primário recebe o tamanho do primário do \texttt{model}; sem
\texttt{maxsize}, o arquivo cresce até ocupar o disco. O
\texttt{create database} também não pode rodar dentro de uma transação
explícita.

\begin{lstlisting}
create database rascunho   -- tamanhos e caminhos padrão
\end{lstlisting}

\subsection{\texttt{alter database}}

Todo \texttt{alter} pressupõe um banco em uso: é comando de manutenção,
corretiva ou por mudança de regra de negócio. Em desenvolvimento, apagar
e recriar costuma ser mais simples. Remover arquivo ou reduzir tamanho
pode perder dados, então a recomendação da aula é fazer backup antes.

\begin{lstlisting}[caption={Alterações comuns nos arquivos de um banco.}]
-- acrescentar um arquivo de dados secundário
alter database biblioteca
add file
   (name = biblioteca_dados2,
    filename = 'C:\dados\biblioteca2.ndf',
    size = 5MB)

-- aumentar o tamanho de um arquivo
alter database biblioteca
modify file (name = biblioteca_dados, size = 20MB)

-- remover um arquivo (ele precisa estar vazio)
alter database biblioteca
remove file biblioteca_dados2
\end{lstlisting}

\subsection{\texttt{drop database}}

Nas palavras da professora, o \texttt{drop} é o comando mais simples e o
mais perigoso, para qualquer objeto.

\begin{lstlisting}
drop database biblioteca
drop database banco1, banco2   -- vários de uma vez
\end{lstlisting}

\subsection{Pelo SSMS}

\begin{itemize}
	\item criar: botão direito em \emph{Databases} $\rightarrow$ \emph{New
	Database};
	\item alterar: botão direito no banco $\rightarrow$ \emph{Properties};
	\item excluir: botão direito no banco $\rightarrow$ \emph{Delete}.
\end{itemize}

\setchapterpreamble[u]{\margintoc}
\chapter{T-SQL: Definição de Dados (DDL)}
\labch{cet615-ddl}

Este capítulo cobre a terceira apresentação da primeira
aula~\cite{BaiocoCET615}: tipos de dados, tabelas, restrições de
integridade e índices. É o projeto físico na prática: transformar o
esquema lógico em um script DDL e executá-lo no SQL Server. Os exemplos
usam um banco de biblioteca, criado para estas anotações.

\section{Tipos de dados}

\begin{table}[h]
	\begin{tabular}{lll}
		\toprule
		Categoria & Tipos & Tamanho \\
		\midrule
		Caractere & \texttt{char(n)}, \texttt{varchar(n)} & até $n$ bytes \\
		Numérico exato & \texttt{decimal(p,e)}, \texttt{numeric(p,e)} & depende de $p$ \\
		Numérico aproximado & \texttt{float}, \texttt{real} & 8 e 4 bytes \\
		Inteiro & \texttt{int}, \texttt{smallint}, \texttt{tinyint} & 4, 2 e 1 byte \\
		Monetário & \texttt{money}, \texttt{smallmoney} & 8 e 4 bytes \\
		Data e hora & \texttt{datetime}, \texttt{smalldatetime} & 8 e 4 bytes \\
		Binário & \texttt{binary(n)}, \texttt{varbinary(n)} & $n$ bytes \\
		Outros & \texttt{bit} & 1 bit \\
		\bottomrule
	\end{tabular}
	\caption{Tipos básicos do SQL Server apresentados em aula.}
\end{table}

Em \texttt{numeric(p,e)}, $p$ é o total de dígitos e $e$, quantos ficam
depois da vírgula: \texttt{numeric(6,2)} vai até 9999,99. Diferente de
\texttt{float} e \texttt{real}, os tipos exatos não têm erro de
arredondamento, e por isso servem para valores monetários.

\subsection{\texttt{char} ou \texttt{varchar}?}

Foi o ponto mais longo da aula sobre tipos\sidecite{BaiocoCET615Aulas}.
Quando todas as colunas têm tamanho fixo, o SGBD sabe exatamente onde
começa cada coluna do registro. Basta uma coluna \texttt{varchar} para o
registro passar a ter tamanho variável: o tamanho de cada valor precisa
ser gravado e lido antes do próprio valor, e cada \texttt{varchar} a mais
soma esse custo. A regra da professora é usar \texttt{varchar} só quando
o ganho supera o custo: quando o tamanho varia muito (de 1 a 100
caracteres, por exemplo) ou quando a coluna é opcional e quase sempre
vazia. Nas palavras dela, tudo o que é muito flexível tem um custo; se
\texttt{varchar} fosse sempre melhor, \texttt{char} não existiria.

Ela também lembrou que \texttt{money} e os tipos de data e hora dependem
da configuração de idioma do servidor.

\subsection*{Complemento: tipos obsoletos e tipos novos}

A tabela dos slides inclui tipos que a Microsoft marca como obsoletos:

\begin{itemize}
	\item \texttt{text}, \texttt{ntext} e \texttt{image} serão removidos em
	versão futura; a documentação manda usar \texttt{varchar(max)},
	\texttt{nvarchar(max)} e \texttt{varbinary(max)}\sidecite{MSTextImage};
	\item \texttt{timestamp} não guarda data nem hora: é sinônimo de
	\texttt{rowversion}, um contador binário que muda a cada alteração da
	linha. A sintaxe \texttt{timestamp} também está obsoleta, e para
	registrar data e hora a recomendação é
	\texttt{datetime2}\sidecite{MSRowversion}.
\end{itemize}

Desde o SQL Server 2008 existem também \texttt{date} (só data),
\texttt{time} (só hora), \texttt{datetime2} (maior precisão) e
\texttt{datetimeoffset} (com fuso). O esquema da primeira avaliação usa
\texttt{date}.

\section{Tabelas}

\subsection{\texttt{create table}}

\begin{lstlisting}[caption={Tabela com nulidade explícita e valor padrão.}]
use biblioteca
go

create table leitor (
codleitor int not null,
nome char(40) not null,
email varchar(80) null,       -- opcional e de tamanho muito variável
cidade char(30) not null default 'Campinas',
primary key (codleitor)
)
\end{lstlisting}

A tabela é criada no banco em uso (\texttt{use}). Cada coluna pode ser
\texttt{null} ou \texttt{not null} e pode ter um \texttt{default}.

O padrão de nulidade, quando nada é declarado, depende da configuração
do banco no servidor. Por isso a professora recomenda declarar
\texttt{null} ou \texttt{not null} em \emph{toda} coluna: o script pode
rodar em outro servidor, com outra configuração. Nomes de coluna não
levam espaço.

\subsection{\texttt{alter table}}

\begin{lstlisting}[caption={Alterações de estrutura.}]
alter table leitor alter column nome char(60) not null  -- muda o tipo
alter table leitor add telefone char(11) null             -- nova coluna
alter table leitor add default 'SP' for uf                -- novo default
alter table leitor drop column telefone                   -- remove coluna

exec sp_rename 'leitor', 'leitores'           -- renomeia tabela
exec sp_rename 'leitores.cidade', 'municipio' -- renomeia coluna
\end{lstlisting}

Regras e alertas:

\begin{itemize}
	\item a coluna nova entra no fim da tabela e, segundo a aula, precisa
	aceitar nulo: as linhas que já existem não teriam valor para ela. O SQL
	Server também aceita uma coluna \texttt{not null} nova se ela tiver
	\texttt{default}; as linhas existentes recebem o valor padrão
	(\texttt{alter table livro add qtd\_exemplares int not null default
	0});
	\item uma coluna só sai se nada depender dela (restrição, índice,
	\texttt{unique}, \texttt{default});
	\item alterações seguras são as que não prejudicam dados existentes:
	aumentar o tamanho de um \texttt{char}, acrescentar coluna ou
	\texttt{default};
	\item renomear tabela ou coluna de um banco em uso pode quebrar as
	aplicações que usam o nome antigo. Vale pensar nos nomes antes;
	\item como no \refch{cet615-sqlserver}, \texttt{alter} é manutenção de
	banco em uso. Em desenvolvimento, a professora prefere apagar e
	recriar.
\end{itemize}

\subsection{\texttt{drop table}}

\begin{lstlisting}
drop table emprestimo
drop table emprestimo, exemplar, livro   -- várias de uma vez
\end{lstlisting}

Uma tabela referenciada por chave estrangeira não pode ser apagada antes
da tabela que a referencia. A ordem de remoção é a inversa da ordem de
criação.

\section{Restrições de integridade}

As restrições de entidade, de chave e referencial do
\refch{cet615-conceitos} viram cláusulas do \texttt{create table}.

\begin{description}
	\item[\texttt{primary key}:]\index{primary key@\texttt{primary key}} a
	coluna é sempre \texttt{not null}, e o SQL Server cria um índice
	\emph{clusterizado} para ela (o padrão pode ser trocado para não
	clusterizado).
	\item[\texttt{unique}:] chave secundária. Pode haver várias por tabela,
	a coluna aceita nulo, e o SQL Server cria um índice \emph{não
	clusterizado} único para ela.
	\item[\texttt{foreign key}:] a tabela referenciada precisa existir
	antes. O SQL Server \emph{não} cria índice para chave estrangeira.
\end{description}

\begin{lstlisting}[caption={Esquema de biblioteca com as três restrições.}]
create table autor (
codautor int not null,
nome char(40) not null,
primary key (codautor)
)

create table livro (
codlivro int not null,
titulo char(60) not null,
isbn numeric(13,0) not null,
codautor int not null,
primary key (codlivro),
unique (isbn),
foreign key (codautor) references autor
)

create index ix_livro_autor on livro (codautor)  -- índice da FK

create table emprestimo (
codemprestimo int not null,
codlivro int not null,
codleitor int not null,
data date not null,
primary key (codemprestimo),
foreign key (codlivro) references livro,
foreign key (codleitor) references leitor
)

create index ix_emprestimo_livro on emprestimo (codlivro)
create index ix_emprestimo_leitor on emprestimo (codleitor)
\end{lstlisting}

Pontos destacados em aula:

\begin{itemize}
	\item Sem chave primária, chave secundária e, principalmente, chave
	estrangeira, o banco deixa de ser relacional. Nas palavras dela, para
	criar tabela sem restrição, melhor usar um banco não relacional.
	\item \texttt{references autor}, sem coluna, aponta para a chave
	primária, como manda o modelo, e falha se a primária tiver sido
	esquecida. Indicar a coluna (\texttt{references autor (codautor)})
	permite apontar para qualquer coluna \texttt{unique}.
	\item O índice da chave estrangeira é obrigatório na prática: o SQL
	Server não o cria. A explicação da professora é que o índice da
	chave estrangeira normalmente não é único, já que quase toda chave
	estrangeira vem de um relacionamento 1:N ou N:M. Num 1:1, o lado 1
	precisa de índice único ou \texttt{unique} na chave estrangeira.
	\item A ordem do script importa: primeiro as tabelas referenciadas,
	depois as que referenciam.
	\item A professora declara as restrições no fim da tabela. Há uma
	armadilha em declará-las na linha da coluna quando a chave é composta:
	duas \texttt{primary key} dão erro, mas \texttt{unique} e
	\texttt{foreign key} não dão erro e criam uma restrição separada por
	coluna, em vez de uma restrição composta.
\end{itemize}

\subsection{Nomear, remover e acrescentar restrições}

Dar nome à restrição é opcional; sem nome, o SQL Server gera um, que
\texttt{sp\_help tabela} mostra. O nome é necessário para remover:

\begin{lstlisting}
exec sp_help emprestimo                  -- descobre o nome
alter table emprestimo drop constraint FK__emprestim__codli__1234

alter table leitor
add constraint uq_leitor_email unique (email)  -- nome explícito
\end{lstlisting}

Restrições normalmente nascem no \texttt{create table}. Criar uma
chave primária depois que o banco está no ar não faz sentido, porque
nulos e repetidos já podem ter entrado. O \texttt{alter table ... add}
serve para corrigir uma chave secundária esquecida (correndo o risco de
já haver repetidos) ou para criar as tabelas em qualquer ordem e
declarar todas as restrições no fim. É assim que o SSMS gera o script de
um banco existente.

\subsection*{Complemento: ações referenciais e \texttt{check}}

\textcite[seção 4.2]{Elmasri2011} trata dois recursos que a aula não
cobre:

\begin{description}
	\item[Ações referenciais.] Por padrão, o SGBD rejeita a operação que
	violaria a integridade referencial: não se apaga um autor que tem
	livros. A chave estrangeira pode declarar outra ação, separada para
	\texttt{on delete} e \texttt{on update}: \texttt{cascade} (propaga a
	exclusão ou a nova chave), \texttt{set null} ou \texttt{set default}.
	O livro sugere \texttt{cascade} para tabelas que mapeiam
	relacionamentos, atributos multivalorados e entidades fracas. No SQL
	Server, o padrão se chama \texttt{no action}.
	\item[\texttt{check}.] Restrição sobre os valores de uma linha, como
	\texttt{check (preco > 0)}. É a forma declarativa das regras
	simples; regras que envolvem outras tabelas ficam para os gatilhos
	(\refch{cet615-gatilhos}).
\end{description}

\begin{lstlisting}
create table exemplar (
codexemplar int not null,
codlivro int not null,
estado char(1) not null check (estado in ('N', 'U', 'D')),
primary key (codexemplar),
foreign key (codlivro) references livro on delete cascade
)
\end{lstlisting}

\section{Índices}

Um \textbf{índice}\index{índice} é uma estrutura que acelera o acesso às
linhas a partir dos valores de uma ou mais colunas. Ele guarda uma
cópia ordenada dos valores da coluna com ponteiros para as linhas.
Ajuda consultas e também atualizações e exclusões, que precisam
encontrar a linha antes de mudá-la. Em troca, ocupa disco e precisa ser
atualizado a cada \texttt{insert}, \texttt{update} e \texttt{delete}.
Muitos índices inúteis pioram as atualizações sem ganho nas consultas.

O uso é transparente: ninguém pede ao SGBD que use o índice. O
\textbf{otimizador de consultas} decide, e pode ignorar um índice que
não ajuda.

\subsection{Quando usar}

\begin{itemize}
	\item colunas usadas no \texttt{where} de \texttt{select},
	\texttt{update} e \texttt{delete}, desde que a condição devolva
	\emph{poucas} linhas;
	\item colunas usadas em \texttt{order by};
	\item colunas usadas em junções.
\end{itemize}

E quando não usar: em tabelas pequenas e em colunas com poucos valores
distintos. Uma coluna com três valores possíveis devolve, em média, um
terço das linhas por consulta, e a varredura sequencial sai mais barata.
A regra empírica da professora é que, acima de 20\% a 25\% das linhas, a
varredura compensa; tabela pequena, para ela, tem menos de mil linhas.

Resumo da aula: chave primária, chave secundária e chave estrangeira se
indexam ``sem pensar''. O que sobra são índices não únicos e não
clusterizados em colunas buscadas com frequência e com poucos repetidos,
como nome de cliente ou título de livro. Uma coluna como UF ou CEP
repete demais para valer a pena.

\subsection{Clusterizado e não clusterizado}

\begin{description}
	\item[Clusterizado (\emph{clustered}):]\index{índice clusterizado} as
	folhas da árvore B são as próprias páginas de dados, na ordem do
	índice. Por isso só pode haver um por tabela: os dados não podem estar
	fisicamente ordenados por duas colunas ao mesmo tempo. Convém criá-lo
	antes dos outros, porque sua criação reordena a tabela e reconstrói os
	demais índices. A escolha deve ser a coluna mais usada em
	\texttt{where}, junções, \texttt{order by} e \texttt{group by}; na
	falta de melhor, a chave primária.
	\item[Não clusterizado (\emph{nonclustered}):] as folhas guardam
	ponteiros para as linhas, que seguem a ordem do clusterizado. Pode
	haver vários por tabela, e é o tipo padrão do \texttt{create index}.
\end{description}

\begin{figure*}[h]
\centering
\begin{tikzpicture}[
	pag/.style={draw=bdroxo, fill=bdroxoclaro, minimum width=14mm, minimum height=6mm, font=\scriptsize},
	dado/.style={draw=bdazul, fill=bdazulclaro, minimum width=14mm, minimum height=6mm, font=\scriptsize},
	lig/.style={draw=bdcinza}]
	% clusterizado
	\node[pag] (r) at (2.25,0) {1 | 40 | 80};
	\node[pag] (i1) at (0,-1.1) {1 | 20};
	\node[pag] (i2) at (2.25,-1.1) {40 | 60};
	\node[pag] (i3) at (4.5,-1.1) {80 | 95};
	\foreach \n in {i1,i2,i3} \draw[lig] (r.south) -- (\n.north);
	\node[dado] (d1) at (0,-2.3) {linhas 1--39};
	\node[dado] (d2) at (2.25,-2.3) {linhas 40--79};
	\node[dado] (d3) at (4.5,-2.3) {linhas 80--};
	\foreach \a/\b in {i1/d1,i2/d2,i3/d3} \draw[lig] (\a.south) -- (\b.north);
	\node[rotulo, align=center] at (2.25,-3.1) {clusterizado: as folhas \emph{são}\\as páginas de dados, em ordem};
	% não clusterizado
	\begin{scope}[xshift=8cm]
	\node[pag] (s) at (2.25,0) {Ana | Lia};
	\node[pag] (j1) at (0.9,-1.1) {Ana | Bia};
	\node[pag] (j2) at (3.6,-1.1) {Lia | Rui};
	\draw[lig] (s.south) -- (j1.north);
	\draw[lig] (s.south) -- (j2.north);
	\node[dado] (e1) at (0,-2.3) {pág. 7};
	\node[dado] (e2) at (1.5,-2.3) {pág. 2};
	\node[dado] (e3) at (3,-2.3) {pág. 9};
	\node[dado] (e4) at (4.5,-2.3) {pág. 4};
	\draw[seta] (j1.south) -- (e1.north);
	\draw[seta] (j1.south) -- (e3.north);
	\draw[seta] (j2.south) -- (e2.north);
	\draw[seta] (j2.south) -- (e4.north);
	\node[rotulo, align=center] at (2.25,-3.1) {não clusterizado: as folhas \emph{apontam}\\para linhas em outra ordem};
	\end{scope}
\end{tikzpicture}
\caption{Índice clusterizado (pela chave primária) e não clusterizado
(pelo nome). Só pode haver um clusterizado, porque os dados só têm uma
ordem física.}
\end{figure*}

\subsection{Único e composto}

\begin{description}
	\item[Único:] os valores não se repetem; a criação falha se já houver
	repetidos. Índice único em coluna \texttt{not null} equivale, na
	prática, a uma chave primária.
	\item[Composto:] formado por várias colunas, em ordem. Funciona como
	ordenar pela primeira coluna e, nos empates, pela segunda. Por isso
	(\texttt{uf}, \texttt{cidade}) é diferente de (\texttt{cidade},
	\texttt{uf}), e o otimizador usa o índice mesmo quando só a primeira
	coluna aparece na busca. Se for único, o que não se repete é a
	combinação.
\end{description}

Uma consequência aparece no script da aula: a chave primária composta
de uma tabela de itens (\texttt{num\_pedido}, \texttt{codigo\_p}) já
gera um índice que atende buscas só por \texttt{num\_pedido}. Só a
segunda chave estrangeira precisa de índice próprio. Pelo mesmo motivo,
a professora reconheceu que o exemplo de índice composto do slide, sobre
as mesmas colunas da chave primária, só serve para mostrar a sintaxe.

\subsection{\texttt{create index} e \texttt{drop index}}

\begin{lstlisting}
create [unique] [clustered | nonclustered] index nome
on tabela (coluna1 [, coluna2, ...])

create index ix_livro_titulo on livro (titulo)
create unique index ix_leitor_cpf on leitor (cpf)
create index ix_exemplar_livro_estado on exemplar (codlivro, estado)

drop index livro.ix_livro_titulo
drop index livro.ix_livro_titulo, leitor.ix_leitor_cpf
\end{lstlisting}

O \texttt{drop index} usa \texttt{tabela.indice} porque o mesmo nome de
índice pode existir em tabelas diferentes. Um índice único em
\texttt{cpf} tem o mesmo efeito de declarar \texttt{unique (cpf)} na
tabela.

\subsection*{Complemento: páginas, limites e a teoria dos índices}

A documentação corrige dois números ditos em aula:

\begin{itemize}
	\item o limite de índices não clusterizados por tabela é 999 desde o
	SQL Server 2008; os 249 do slide eram o limite das versões
	anteriores\sidecite{MSCapacity};
	\item a página do SQL Server tem tamanho fixo de 8 KB e não é
	configurável; oito páginas contíguas formam uma \emph{extensão} de
	64 KB, e a leitura e a gravação em disco são feitas por página
	inteira\sidecite{MSPages}. Uma linha comporta até 8.060 bytes. O que
	o DBA ajusta é o \emph{fill factor}, a fração de cada página de índice
	preenchida na criação.
\end{itemize}

Na teoria, \textcite[cap. 18]{Elmasri2011} distingue três índices de um
nível: o \emph{primário}, sobre o campo-chave de um arquivo ordenado por
ele; o de \emph{agrupamento} (\emph{clustering}), sobre um campo não
chave pelo qual o arquivo está ordenado; e o \emph{secundário}, sobre
qualquer outro campo. Um arquivo tem no máximo um índice primário ou de
agrupamento, e quantos secundários quiser. É a mesma restrição do índice
clusterizado do SQL Server. O capítulo também explica as árvores B e
B$^+$, a estrutura dos índices de vários níveis que o SQL Server usa.

\section{Do esquema ao script}

O roteiro da aula para transformar o esquema lógico em DDL:

\begin{enumerate}
	\item criar as tabelas na ordem do esquema lógico, das referenciadas
	para as que referenciam;
	\item em cada tabela, declarar tipos, nulidade, chave primária,
	chaves secundárias e chaves estrangeiras;
	\item criar um índice para cada chave estrangeira que não seja a
	primeira coluna da chave primária;
	\item avaliar índices extras pelas consultas previstas;
	\item conferir o resultado no diagrama do SSMS (\emph{Database
	Diagrams} $\rightarrow$ \emph{New Database Diagram}), que desenha o
	esquema físico a partir do catálogo;
	\item para desfazer, apagar as tabelas na ordem inversa.
\end{enumerate}

\setchapterpreamble[u]{\margintoc}
\chapter{T-SQL: Manipulação de Dados (DML)}
\labch{cet615-dml}

Este capítulo cobre a primeira metade da segunda
aula~\cite{BaiocoCET615}: inserir, alterar e excluir dados; consultar
com filtros, agrupamentos, junções e subconsultas; e agrupar operações em
transações. Os exemplos de consulta da aula usam o banco PUBS, de
exemplo da Microsoft. Aqui os exemplos usam o banco de biblioteca do
\refch{cet615-ddl}.

\section{Inserir, alterar e excluir}

\subsection{\texttt{insert}}

\begin{lstlisting}
-- sem lista de colunas: todos os valores, na ordem da tabela
insert into leitor values (1, 'Ana Lima', null, 'Campinas')

-- com lista de colunas: as omitidas recebem nulo ou o default
insert into leitor (codleitor, nome) values (2, 'Bruno Reis')

-- várias linhas no mesmo comando
insert into leitor (codleitor, nome)
values (3, 'Carla Dias'), (4, 'Davi Rocha')

-- resultado de uma consulta
insert into leitor_vip (codleitor, nome)
select codleitor, nome from leitor where cidade = 'Campinas'
\end{lstlisting}

Texto e data vão entre aspas; número, não. \texttt{null} é um valor
especial que significa ``não informado'': não é branco nem zero. Gravar
\texttt{''} ou \texttt{0} grava exatamente isso. No \texttt{insert} sem
lista de colunas, o nulo precisa ser escrito na posição da coluna.

Duas observações da aula\sidecite{BaiocoCET615Aulas}: o \texttt{insert}
com várias linhas tem o mesmo resultado de vários \texttt{insert}, mas
roda como uma única transação; e o \texttt{insert ... select} serve, por
exemplo, para montar uma tabela desnormalizada para consulta rápida ou
para uma ferramenta de BI.

\subsection{\texttt{update} e \texttt{delete}}

\begin{lstlisting}
update livro set titulo = 'Dom Casmurro' where codlivro = 10
update exemplar set estado = 'U', codlivro = 11 where codexemplar = 5
update livro set preco = preco * 1.1        -- todas as linhas!

delete from emprestimo where codemprestimo = 100
delete emprestimo                            -- o from é opcional
\end{lstlisting}

O \texttt{where} é ``quase obrigatório'': sem ele, o comando atinge todas
as linhas. Filtrar pela chave garante uma linha só; filtrar por outra
coluna pode atingir várias. Um \texttt{update} ou \texttt{delete} que não
encontra nada não dá erro: só informa ``0 linhas afetadas''. Esse
comportamento volta nas transações (\refsec{cet615-transacoes-sqlserver}).

\subsection{\texttt{truncate table}}

Para esvaziar uma tabela grande, \texttt{delete} sem \texttt{where} pode
estourar o log de transações, que é bem menor que o banco. O comando
indicado é \texttt{truncate table emprestimo}.

\subsection*{Complemento: o que o \texttt{truncate} registra}

Em aula, a professora disse que o \texttt{truncate} não usa o log; o
slide diz que ele usa menos log. A documentação\sidecite{MSTruncate}
confirma o slide: o \texttt{delete} grava no log uma entrada por linha,
enquanto o \texttt{truncate} desaloca as páginas da tabela e registra só
essas desalocações. Outras diferenças documentadas:

\begin{itemize}
	\item o \texttt{truncate} pode ser desfeito com \texttt{rollback}
	dentro de uma transação;
	\item não funciona em tabela referenciada por chave estrangeira;
	\item não dispara gatilhos, porque não registra as linhas
	excluídas;
	\item reinicia a contagem de colunas \texttt{identity}.
\end{itemize}

\section{Consultas}

\subsection{Estrutura do \texttt{select}}

\begin{lstlisting}
select [distinct] colunas      -- o que mostrar
from tabelas                   -- de onde
[where condição]               -- quais linhas
[group by colunas]             -- como agrupar
[having condição]              -- quais grupos
[order by colunas [asc|desc]]  -- em que ordem
\end{lstlisting}

Antes de consultar, escolha o banco com \texttt{use}. A professora
lembrou que conhecer o esquema (chaves e relacionamentos) é o primeiro
passo para escrever consultas, principalmente com junção.

\subsection{SQL e álgebra relacional}

Toda consulta SQL é traduzida pelo SGBD em operações de álgebra
relacional. Na compilação, o SGBD considera árvores de álgebra
equivalentes e escolhe a mais barata; a otimização não depende só de
índices e estatísticas. A correspondência mostrada em aula:

\begin{description}
	\item[Lista do \texttt{select} = projeção ($\pi$):] corte vertical,
	escolhe colunas.
	\item[\texttt{where} = seleção ($\sigma$):] corte horizontal, escolhe
	linhas.
	\item[Junção no \texttt{from}:] produto cartesiano seguido de uma
	condição, mas executado por um algoritmo próprio, mais eficiente.
\end{description}

Na álgebra, o resultado de uma consulta é uma relação e, portanto, não
tem tuplas repetidas. Os SGBDs não eliminam repetições por padrão (seria
caro), por isso o \texttt{distinct} precisa ser pedido.
\textcite[cap. 6]{Elmasri2011} cobre a álgebra relacional, que a
professora recomendou revisar.

\subsection{Filtros}

\begin{lstlisting}
select titulo from livro where preco > 50
select titulo from livro where preco between 20 and 50
select nome from leitor where cidade in ('Campinas', 'Jundiaí')
select nome from leitor where cidade not in ('Campinas')
select nome from leitor where nome like 'A%'      -- começa com A
select nome from leitor where nome like '_na%'    -- 2.a e 3.a letras "na"
select nome from leitor where nome like '[A-C]%'  -- começa com A, B ou C
\end{lstlisting}

\texttt{between}, \texttt{in} e \texttt{like} com colchetes são
atalhos: tudo pode ser escrito com \texttt{>=}, \texttt{<=}, \texttt{=}
e \texttt{or}, só que mais longo. No \texttt{like}, \texttt{\%} vale
qualquer sequência (inclusive vazia) e \texttt{\_} vale exatamente um
caractere naquela posição, o que fixa o tamanho do texto.

A professora frisou a \textbf{precedência}: operadores relacionais vêm
antes dos lógicos, e \texttt{and} vem antes de \texttt{or}.

\begin{lstlisting}
-- errado: o filtro de preço só vale para 'romance'
select titulo from livro
where genero = 'poesia' or genero = 'romance' and preco > 50

-- certo
select titulo from livro
where (genero = 'poesia' or genero = 'romance') and preco > 50
\end{lstlisting}

Parênteses a mais não fazem mal; a menos, podem dar resultado errado.

\subsection{Nomes, cálculos e nulos}

\begin{lstlisting}
select titulo, preco * 0.9 as preco_desconto from livro
select titulo, preco * 0.9 "preço com desconto" from livro
select titulo from livro where preco is null
\end{lstlisting}

O \texttt{as} é opcional, e as aspas só são necessárias quando o nome tem
espaço. O apelido é útil sobretudo em colunas calculadas, que sairiam sem
nome. No SQL Server, o apelido não vale no \texttt{where} (é preciso
repetir a expressão), mas vale no \texttt{order by}.

Para testar nulo, use \texttt{is null} e \texttt{is not null}: a
comparação \texttt{preco = null} não traz linha nenhuma.

\subsection*{Complemento: a lógica de três valores}

\textcite[seção 5.1.1]{Elmasri2011} explica por que \texttt{= null} não
funciona. Um nulo pode significar valor desconhecido, valor omitido de
propósito ou atributo que não se aplica, e a SQL não distingue os três
casos. Toda comparação com nulo dá \textsc{desconhecido}, e não
verdadeiro ou falso. A SQL usa então uma lógica de três valores:

\begin{itemize}
	\item \textsc{falso} \texttt{and} qualquer coisa dá \textsc{falso};
	\textsc{verdadeiro} \texttt{and} \textsc{desconhecido} dá
	\textsc{desconhecido};
	\item \textsc{verdadeiro} \texttt{or} qualquer coisa dá
	\textsc{verdadeiro}; \textsc{falso} \texttt{or} \textsc{desconhecido}
	dá \textsc{desconhecido};
	\item \texttt{not} \textsc{desconhecido} continua
	\textsc{desconhecido}.
\end{itemize}

O \texttt{where} só devolve as linhas em que a condição é
\textsc{verdadeiro}. Por isso \texttt{not in} tem uma armadilha: se a
lista (ou a subconsulta) contém um nulo, a condição nunca é verdadeira e
a consulta volta vazia.

\subsection{Ordenação}

\begin{lstlisting}
select genero, titulo, preco from livro order by genero, preco desc
\end{lstlisting}

O padrão é \texttt{asc}. No SQL Server, os nulos vêm primeiro na ordem
ascendente e por último na descendente. Ordenar pela segunda coluna só
muda algo se a primeira tiver valores repetidos; ordenar por chave e
depois por outra coluna gasta recurso à toa.

\subsection{Funções de agregação, \texttt{group by} e \texttt{having}}

\begin{lstlisting}
select count(*) from livro              -- linhas
select count(preco) from livro          -- linhas com preço não nulo
select avg(preco), min(preco), max(preco), sum(preco) from livro
select avg(distinct preco) from livro

select genero, avg(preco) as media
from livro
group by genero
having avg(preco) > 40
\end{lstlisting}

\begin{itemize}
	\item \texttt{count(*)} conta linhas; \texttt{count(coluna)} e as
	outras funções ignoram nulos.
	\item Sem \texttt{group by}, \texttt{select genero, avg(preco)} dá
	erro: a média devolve um valor e \texttt{genero}, vários. Toda coluna
	fora de função de agregação precisa estar no \texttt{group by}.
	\item Agrupar por chave não faz sentido: cada grupo teria uma linha só.
	\item O \texttt{where} filtra linhas \emph{antes} do agrupamento e não
	enxerga os grupos; para filtrar grupos existe o \texttt{having}. A
	condição do \texttt{having} pode usar uma função diferente da que está
	no \texttt{select}.
\end{itemize}

\section{Junções}

O \textbf{produto cartesiano} (\texttt{cross join}) combina cada linha
de uma tabela com todas as da outra: 18 títulos e 8 editoras dão 144
linhas. Sozinho, quase nunca tem uso. A \textbf{junção}\index{junção}
é o produto seguido de uma condição, geralmente igualdade entre chave
estrangeira e chave primária. O algoritmo, porém, não monta o produto
inteiro para filtrar depois: testa a condição a cada combinação e
descarta na hora o que não serve.

\begin{lstlisting}[caption={Junção interna com duas e três tabelas.}]
select l.titulo, a.nome
from livro l inner join autor a
     on l.codautor = a.codautor
where l.preco > 50

select le.nome, l.titulo, e.data
from leitor le inner join emprestimo e
     on le.codleitor = e.codleitor
     inner join livro l
     on e.codlivro = l.codlivro
\end{lstlisting}

\begin{itemize}
	\item A condição de junção fica no \texttt{on}; os filtros ficam no
	\texttt{where}.
	\item Junta-se de duas em duas tabelas, seguindo os relacionamentos.
	Uma tabela já incluída fica visível para os \texttt{on} seguintes.
	\item Uma tabela não tem colunas com o mesmo nome, mas o resultado de
	uma junção pode ter. Sem o prefixo (\texttt{l.codautor}), dá erro de
	ambiguidade. O apelido da tabela deve ser curto.
	\item \texttt{join} sozinho é \texttt{inner join}.
	\item A sintaxe antiga, com vírgula no \texttt{from} e a condição no
	\texttt{where}, ainda funciona, mas mistura junção com filtro e fica
	confusa com três tabelas ou mais. A professora pede a notação
	\texttt{inner join ... on}.
\end{itemize}

\subsection{Junções externas}

\begin{lstlisting}
-- todos os leitores, inclusive os sem empréstimo
select le.nome, e.codemprestimo
from leitor le left join emprestimo e
     on le.codleitor = e.codleitor

-- só os leitores sem empréstimo
select le.nome
from leitor le left join emprestimo e
     on le.codleitor = e.codleitor
where e.codemprestimo is null
\end{lstlisting}

\begin{figure}[h]
\centering
\begin{tikzpicture}[
	tab/.style={draw=bdcinza, font=\scriptsize, minimum width=17mm, minimum height=5mm, inner sep=1pt},
	ok/.style={tab, fill=bdazulclaro},
	nul/.style={tab, fill=white, text=bdcinza}]
	\node[font=\scriptsize\bfseries] at (0,0.6) {leitor};
	\node[tab] (a1) at (0,0) {1 Ana};
	\node[tab, below=0pt of a1] (a2) {2 Bruno};
	\node[tab, below=0pt of a2] (a3) {3 Carla};
	\node[font=\scriptsize\bfseries] at (2.3,0.6) {empréstimo};
	\node[tab] (b1) at (2.3,0) {100 \textrightarrow{} 1};
	\node[tab, below=0pt of b1] (b2) {101 \textrightarrow{} 3};
	\node[tab, below=0pt of b2] (b3) {102 \textrightarrow{} 9};
	% inner
	\node[font=\scriptsize\bfseries] at (5.2,0.6) {inner join};
	\node[ok] (i1) at (5.2,0) {Ana | 100};
	\node[ok, below=0pt of i1] (i2) {Carla | 101};
	% left
	\node[font=\scriptsize\bfseries] at (7.5,0.6) {left join};
	\node[ok] (l1) at (7.5,0) {Ana | 100};
	\node[nul, below=0pt of l1] (l2) {Bruno | null};
	\node[ok, below=0pt of l2] (l3) {Carla | 101};
\end{tikzpicture}
\caption{Junção interna e externa à esquerda. Bruno não tem empréstimo e
só aparece no \texttt{left join}, com nulo. O empréstimo 102 aponta para
um leitor inexistente, situação que a chave estrangeira impediria; ele só
apareceria num \texttt{right} ou \texttt{full join}.}
\end{figure}

O \texttt{left join} mantém todas as linhas da tabela da esquerda e
preenche com nulos o que não tem correspondência; o \texttt{right join}
faz o mesmo com a da direita, e o \texttt{full join} com as duas. A
palavra \texttt{outer} é opcional. Para achar as linhas sem
correspondência, testa-se \texttt{is null} numa coluna da tabela do
outro lado, que nem precisa estar no \texttt{select}.

\subsection{O erro que conta na nota}

O alerta com mais peso de avaliação da aula: \emph{não juntar tabela
desnecessária}. Junção é cara. Se o código que já está na tabela de
relacionamento basta, a consulta para ali. Numa consulta otimizada,
juntar uma tabela que não contribui com nenhuma coluna nem filtro conta
como erro, mesmo que o resultado saia certo.

\begin{lstlisting}
-- desnecessário: autor não contribui com nada
select l.titulo, l.codautor
from livro l inner join autor a on l.codautor = a.codautor

-- suficiente
select titulo, codautor from livro
\end{lstlisting}

\section{Subconsultas}

\begin{lstlisting}
-- livros da editora "Nova Lua", sem junção
select titulo from livro
where codeditora = (select codeditora from editora
                    where nome = 'Nova Lua')

-- livros de autores com sobrenome Ramos (pode haver vários)
select titulo from livro
where codautor in (select codautor from autor
                   where sobrenome = 'Ramos')
\end{lstlisting}

A posição da professora é clara:

\begin{itemize}
	\item a subconsulta que vale a pena é a do \texttt{where} que
	\emph{substitui} uma junção: são dois passos simples, sem produto
	cartesiano;
	\item use \texttt{=} só quando tiver certeza de que a subconsulta
	devolve um único valor; senão, \texttt{in};
	\item a consulta externa não enxerga colunas das tabelas internas. Se
	essas colunas são necessárias, a solução é a junção;
	\item subconsultas no \texttt{select} ou no \texttt{from} precisam ser
	muito bem pensadas. Na opinião dela, perdem para a junção na maioria
	dos casos, porque cada uma pode carregar sua própria junção. Ela
	disse ver esse uso em toda parte, possivelmente em código gerado por
	IA, e pediu senso crítico. O tema volta na disciplina de otimização.
\end{itemize}

\textcite[seção 5.1]{Elmasri2011} vai além da aula com subconsultas
correlacionadas, \texttt{exists} e \texttt{unique}.

\section{Transações}

Uma \textbf{transação}\index{transação} é um conjunto de operações que o
usuário enxerga como uma só: ou tudo é feito, ou nada. Exemplos da aula:
uma transferência entre contas (débito e crédito); uma venda com baixa
de estoque; o cadastro de uma pessoa física, que grava a parte genérica
e a específica.

\subsection{Estados}

\begin{description}
	\item[Ativa:] desde o \texttt{begin}, executando leituras e escritas.
	\item[Parcialmente efetivada:] depois da última operação, mas antes do
	\texttt{commit}. Ainda pode falhar (uma queda de energia, por
	exemplo).
	\item[Efetivada:] depois do \texttt{commit}.
	\item[Falha:] depois de um erro, com \texttt{rollback} (ou
	\emph{abort}, sinônimos).
	\item[Encerrada:] o destino de todas, com sucesso ou com falha.
\end{description}

\begin{figure}[h]
\centering
\begin{tikzpicture}[estado/.style={caixa, rounded corners=8pt, minimum width=24mm}]
	\node[estado] (at) at (0,0) {Ativa};
	\node[estado] (pe) at (5,0) {Parcialmente\\efetivada};
	\node[estado] (ef) at (5,-2) {Efetivada};
	\node[estado] (fa) at (0,-2) {Falha};
	\node[estado, fill=bdroxo, text=white] (en) at (2.5,-3.6) {Encerrada};
	\draw[seta] (-2.6,0) -- node[rotulo, above] {begin} (at);
	\draw[seta] (at) -- node[rotulo, above] {última operação} (pe);
	\draw[seta] (pe) -- node[rotulo, right] {commit} (ef);
	\draw[seta] (at) -- node[rotulo, left] {erro} (fa);
	\draw[seta] (pe) -- node[rotulo, above, sloped] {falha} (fa);
	\draw[seta] (ef) -- (en);
	\draw[seta] (fa) -- node[rotulo, left] {rollback} (en);
	\draw[seta] (at) edge[loop above, looseness=5] node[rotulo, above] {leituras e escritas} (at);
\end{tikzpicture}
\caption{Estados de uma transação. Com sucesso ou com falha, toda
transação precisa chegar a encerrada, que libera os bloqueios.}
\end{figure}

A professora insistiu que toda transação precisa chegar ao estado
encerrada: uma transação aberta mantém os bloqueios e deixa os outros
usuários esperando. O \texttt{rollback} ``dá a impressão de desfazer'',
mas é sobretudo o aviso de encerramento.

\subsection{Propriedades ACID}

\begin{description}
	\item[Atomicidade:] tudo ou nada.
	\item[Consistência:] a transação leva o banco de um estado consistente
	a outro.
	\item[Isolamento:] cada transação se comporta como se executasse
	sozinha.
	\item[Durabilidade:] o que foi efetivado não se perde, nem com falha.
\end{description}

Na divisão da aula, atomicidade e consistência dependem de quem escreve
a transação: o SGBD obedece, mas só o projetista sabe o que precisa ser
agrupado. Se o débito de uma transferência for esquecido, o SGBD executa
o crédito sem reclamar. Isolamento e durabilidade são responsabilidade
do SGBD, que enxerga tudo o que está rodando.

\textcite[seção 21.3]{Elmasri2011} divide um pouco diferente: o
subsistema de recuperação garante atomicidade e durabilidade, o controle
de concorrência garante isolamento, e a consistência cabe aos
programadores e ao módulo de restrições de integridade. As duas visões
se completam: o SGBD garante o ``tudo ou nada'' daquilo que o
programador colocou dentro da transação.

\subsection{Conflito e bloqueio}

Duas operações \textbf{conflitam} quando pertencem a transações
diferentes, acessam o mesmo item (tabela, página ou linha, conforme a
granularidade) e pelo menos uma delas é escrita. Leituras não conflitam
entre si. Diante de um conflito, o SGBD serializa as operações.

O mecanismo usual é o \textbf{bloqueio}\index{bloqueio}: compartilhado
para leitura e exclusivo para escrita. Bloquear o banco inteiro seria
absurdo; a tabela é fácil de bloquear; a linha gera menos conflitos, mas
custa mais para gerenciar. A página é o meio-termo, e é a granularidade
mais comum. Bloqueios explícitos, escritos pelo programador, existem,
mas o desejável é deixar o SGBD bloquear sozinho.

Um caso comentado em aula: uma escrita precisa de bloqueio exclusivo e
pode esperar para sempre enquanto chegam leituras com bloqueio
compartilhado (\emph{livelock}). A solução é que, quando há uma escrita
na fila, nenhuma nova leitura recebe bloqueio compartilhado.

Regra prática da professora: transação enxuta, só com o que precisa ser
agrupado, e com fim claro em \texttt{commit} ou \texttt{rollback}.

\subsection{Transações no SQL Server}
\labsec{cet615-transacoes-sqlserver}

\begin{itemize}
	\item O modo padrão é \emph{autocommit}: cada comando é uma transação
	e é efetivado sozinho se não der erro.
	\item A transação explícita começa com \texttt{begin tran} (ou
	\texttt{begin transaction}) e termina com \texttt{commit} ou
	\texttt{rollback}. Não existe \texttt{end transaction}.
	\item No modo implícito, cada \texttt{commit} ou \texttt{rollback} já
	abre a próxima transação. A professora desaconselha: exige muito
	domínio e tende a segurar bloqueios por mais tempo.
	\item \texttt{save tran} cria um ponto de retorno parcial. Ela não
	vê sentido nele numa transação bem projetada.
	\item O bloqueio é implícito; a granularidade padrão é a página.
	\item \texttt{@@rowcount} é a variável do sistema com o número de
	linhas afetadas pelo último comando. Ela usa \texttt{> 0} em vez de
	\texttt{= 1} para valer também em \texttt{update} de várias linhas.
	\item Blocos são \texttt{begin ... end}. Não existe \texttt{end if}.
\end{itemize}

A aula não usa \texttt{try...catch}: a lógica é escrita com
\texttt{if}, que todo SGBD tem, para deixar visível o caminho de cada
falha. O padrão de uma transferência:

\begin{lstlisting}[caption={Transferência entre contas com \texttt{@@rowcount}.}]
begin transaction
update conta set saldo = saldo - 100 where numero = 1
if @@rowcount > 0               -- débito feito
begin
   update conta set saldo = saldo + 100 where numero = 2
   if @@rowcount > 0            -- crédito feito
      commit transaction
   else
      rollback transaction
end
else
   rollback transaction
\end{lstlisting}

Todo caminho termina em \texttt{commit} ou \texttt{rollback}, inclusive
o primeiro \texttt{else}: mesmo sem nada gravado, a transação foi aberta
e precisa ser encerrada. Com chaves estrangeiras, a ordem importa:
insere-se primeiro a linha referenciada (cliente antes de pessoa física,
conta antes de cliente--conta). Na demonstração da aula, a segunda
inserção falhou e a primeira não ficou no banco. Sem a transação, ela
teria ficado gravada sozinha.

\subsection*{Complemento: o que o \texttt{rollback} faz por dentro}

Em aula, a professora explicou que as escritas ficam no log e só vão
para o banco no \texttt{commit}; o \texttt{rollback} apenas as descarta.
É a estratégia de atualização adiada, que \textcite[cap. 23]{Elmasri2011}
chama de NO-UNDO/REDO. O livro descreve também a atualização imediata,
em que páginas alteradas podem ir para o disco antes do \texttt{commit}
e o \texttt{rollback} precisa desfazê-las.

O SQL Server segue a segunda linha\sidecite{MSTransactionLog}. Ele usa
\emph{write-ahead logging}: nenhuma página alterada vai para o disco
antes do registro de log correspondente. As páginas são alteradas na
memória (\emph{buffer cache}) e gravadas depois, por exemplo num
\emph{checkpoint}, que pode acontecer antes do \texttt{commit}. Por isso
o log guarda a imagem anterior de cada alteração, e o \texttt{rollback}
a reaplica. Na recuperação depois de uma queda, o SQL Server refaz o
que foi efetivado e não chegou ao disco e desfaz o que ficou sem
\texttt{commit}. O resultado visto pelo usuário é o mesmo da explicação
da aula: o que não foi efetivado não fica no banco.

A documentação reforça também o alerta sobre transações abertas: o log
não pode ser truncado antes da transação ativa mais antiga, então uma
transação esquecida faz o log crescer sem parar.

\setchapterpreamble[u]{\margintoc}
\chapter{T-SQL: Visões}
\labch{cet615-visoes}

Este capítulo cobre a apresentação sobre visões da segunda
aula\sidecite{BaiocoCET615}. No SSMS em português, as visões aparecem
como \emph{Exibições}.

\section{O que é uma visão}

Uma \textbf{visão}\index{visão} (\emph{view}) é uma consulta guardada
com nome, que pode ser usada como se fosse uma tabela. A imagem da aula
é a de uma moldura: através dela se veem só os dados de interesse de uma
ou mais tabelas (ou de outras visões). O SGBD guarda apenas a definição;
os dados continuam nas tabelas. Uma mudança feita nas tabelas aparece
na visão, e vice-versa.

\subsection{Vantagens}

\begin{itemize}
	\item \textbf{Segurança.} A visão pode esconder colunas e linhas. Em
	aula, a professora foi além do slide: em vez de dar ao usuário
	permissão nas três tabelas de uma consulta, dá-se permissão só na
	visão. O acesso passa a ser por funcionalidade, e o usuário nunca toca
	as tabelas.
	\item \textbf{Simplicidade.} Uma consulta complexa e frequente vira um
	\texttt{select * from visao}.
	\item \textbf{Valores calculados e resumos} ficam prontos na visão.
\end{itemize}

\subsection{Desvantagens}

\begin{itemize}
	\item O desempenho é o do \texttt{select} que define a visão. Se a
	visão tem junção, a junção roda mesmo quando a consulta só precisa de
	colunas de uma das tabelas.
	\item Alterar dados por meio de visões tem restrições (veja adiante).
\end{itemize}

A professora corrigiu a leitura do slide: a visão em si não é ruim para
o desempenho. O problema é usá-la quando a junção dela não é necessária.

\section{Criar, usar, alterar e excluir}

\begin{lstlisting}[caption={Visão sobre o banco de biblioteca.}]
create view emprestimos_detalhados (leitor, livro, data_emprestimo)
as
select le.nome, l.titulo, e.data
from leitor le inner join emprestimo e
     on le.codleitor = e.codleitor
     inner join livro l
     on e.codlivro = l.codlivro
go

select * from emprestimos_detalhados
where data_emprestimo >= '2026-11-01'
\end{lstlisting}

\begin{itemize}
	\item A visão é criada no banco em uso, e o \texttt{create view} precisa
	ser o primeiro comando do lote (\refch{cet615-sqlserver}).
	\item Os nomes das colunas vão na lista do \texttt{create view} ou como
	apelidos no \texttt{select}. São obrigatórios quando duas colunas
	teriam o mesmo nome (dois \texttt{nome} numa junção) e em colunas
	calculadas.
	\item Para criar uma visão, basta saber escrever o \texttt{select}. A
	recomendação é testar o \texttt{select} antes de transformá-lo em
	visão.
	\item O uso ideal é \texttt{select * from visao}, com \texttt{where}
	quando preciso. Um critério que muda com o tempo, como um intervalo de
	datas, fica fora da visão, no \texttt{where} de quem a consulta.
\end{itemize}

\begin{lstlisting}
alter view emprestimos_detalhados (leitor, livro, data_emprestimo)
as
select ...   -- o select inteiro precisa ser reescrito

exec sp_helptext emprestimos_detalhados   -- mostra a definição

drop view emprestimos_detalhados
\end{lstlisting}

O \texttt{alter view} substitui a definição e mantém as permissões já
concedidas. O \texttt{drop} seguido de \texttt{create} perde as
permissões. Por isso o \texttt{alter} só faz sentido com o sistema em
produção; em projeto, apagar e recriar resolve.

\subsection{O antipadrão demonstrado em aula}

A professora mostrou de propósito o uso errado. Para listar só os nomes
dos fornecedores, ela consultou a visão de fornecedores e produtos. O
resultado veio com nomes repetidos, sinal de que a junção estava sendo
executada. Um \texttt{distinct} ``resolve'', mas a forma certa é
consultar a tabela \texttt{fornecedor} diretamente. Já o filtro pelo
nome do produto faz sentido na visão, porque a junção é necessária para
chegar a ele.

\begin{lstlisting}
-- antipadrão: junção desnecessária, remendada com distinct
select distinct leitor from emprestimos_detalhados

-- certo
select nome from leitor
\end{lstlisting}

\section{Alterar dados por meio de visões}

Segundo os slides, uma visão só é atualizável se o seu \texttt{select}:

\begin{itemize}
	\item não tem funções de agregação na lista de seleção;
	\item não tem colunas calculadas a partir de outras;
	\item referencia pelo menos uma tabela no \texttt{from}. Uma visão como
	\texttt{select getdate() as hoje} nunca é atualizável.
\end{itemize}

E os comandos têm restrições:

\begin{itemize}
	\item \texttt{insert} e \texttt{update} só podem alterar uma das
	tabelas do \texttt{from};
	\item o \texttt{insert} precisa satisfazer todas as colunas
	\texttt{not null} das tabelas, mesmo as que não aparecem na visão.
	Como a visão é só uma moldura, a coluna obrigatória que ficou de fora
	fica sem valor, e a inserção falha;
	\item \texttt{delete} não é permitido em visão com mais de uma tabela.
\end{itemize}

A conclusão da professora é que \emph{visão serve para consulta}. Para
usar \texttt{insert} ou \texttt{update} por uma visão, o usuário
precisaria conhecer o \texttt{select} dela para saber em que tabela está
mexendo. E incluir na visão todas as colunas obrigatórias, para o
\texttt{insert} funcionar, desfaz o objetivo de segurança. Para escrita,
há saídas melhores: outro usuário com permissão nas tabelas ou um
procedimento armazenado (\refch{cet615-procedimentos}).

\subsection*{Complemento: como o SGBD executa uma visão}

\textcite[seção 5.3.3]{Elmasri2011} descreve duas estratégias de
implementação. Na \textbf{modificação de consulta}, o SGBD reescreve a
consulta sobre a visão como uma consulta sobre as tabelas de base. Na
\textbf{materialização}, ele grava o resultado da visão numa tabela e a
mantém atualizada de forma incremental enquanto ela é usada.

No SQL Server, a visão comum segue a primeira estratégia: a
documentação\sidecite{MSCreateView} diz que uma operação sobre a visão é
validada e traduzida em uma operação sobre as tabelas de base. Só a
\emph{visão indexada} guarda dados. Isso reforça o ponto da professora:
o custo de usar uma visão é o custo do \texttt{select} que ela esconde.

O livro também explica por que a atualização por visão é restrita: numa
visão com junção, uma mesma alteração pode corresponder a várias
alterações diferentes nas tabelas de base, e o SGBD não tem como saber
qual foi a pretendida.

\subsection*{Complemento: opções do \texttt{create view}}

A documentação\sidecite{MSCreateView} traz recursos que a aula não
cobre:

\begin{description}
	\item[\texttt{with check option}:] impede que uma alteração feita pela
	visão produza uma linha que a própria visão deixaria de mostrar.
	\item[\texttt{with schemabinding}:] amarra a visão às tabelas; elas não
	podem ser alteradas de um jeito que afete a visão, nem apagadas antes
	dela. É pré-requisito da visão indexada.
	\item[\texttt{order by}:] não é aceito na definição, a não ser com
	\texttt{top}, e mesmo assim não garante a ordem do resultado. A ordem
	vai na consulta sobre a visão.
	\item[Gatilhos \texttt{instead of}:] tornam atualizável uma visão que
	não seria, porque o gatilho decide o que fazer com cada alteração
	(\refch{cet615-gatilhos}).
\end{description}

A documentação também atualiza um detalhe do slide: o texto da visão,
que o \texttt{sp\_helptext} mostra, fica na visão de catálogo
\texttt{sys.sql\_modules}; \texttt{syscomments} é uma visão de
compatibilidade.

\setchapterpreamble[u]{\margintoc}
\chapter{T-SQL: Procedimentos Armazenados e Transações}
\labch{cet615-procedimentos}

Este capítulo cobre a última apresentação da segunda
aula\sidecite{BaiocoCET615}: procedimentos armazenados e o uso deles
para encapsular transações. É o conteúdo da questão 2 da primeira
avaliação (\refch{cet615-avaliacoes}).

\section{O que é um procedimento armazenado}

Um \textbf{procedimento armazenado}\index{procedimento armazenado}
(\emph{stored procedure}) é um conjunto de comandos SQL compilado e
guardado no servidor, dentro de um banco. A analogia da aula é a de uma
rotina de linguagem de programação: a lógica fica no banco, e a
aplicação só chama procedimentos e visões, com o mínimo de SQL escrito
nela.

Vantagens:

\begin{itemize}
	\item \textbf{Encapsulamento.} Rotinas frequentes ficam disponíveis
	para todas as aplicações, em vez de reescritas em cada uma.
	\item \textbf{Menos erros.} Em tese, os procedimentos já foram
	testados. Em aula, a professora disse esperar que sim na prática, e
	que pelo menos os erros ficam mais fáceis de localizar.
	\item \textbf{Segurança.} O acesso fica limitado ao que o
	procedimento faz, não importa quantas tabelas ele use. Com visões e
	procedimentos, a permissão é dada por funcionalidade, como a interface
	libera ou bloqueia uma opção, e não por tabela.
\end{itemize}

No SSMS em português, os procedimentos ficam em \emph{Programação}
$\rightarrow$ \emph{Procedimentos Armazenados}.

\section{Criar e executar}

\begin{lstlisting}[caption={Busca de livros por parte do título.}]
create procedure busca_livro
@trecho varchar(60)
as
select codlivro, titulo
from livro
where titulo like '%' + @trecho + '%'
go

exec busca_livro 'casmurro'
execute busca_livro @trecho = 'casmurro'   -- passagem por nome
\end{lstlisting}

O parâmetro é \texttt{varchar} de propósito. Um \texttt{char(60)} viria
completado com espaços à direita, e \texttt{'\%' + @trecho + '\%'}
procuraria o trecho seguido de brancos. Os parâmetros que vão para
colunas, porém, devem ter o mesmo tipo da coluna.

A passagem pode ser posicional ou por nome. Depois que a chamada começa
a passar parâmetros por nome, os seguintes também precisam ser por nome.
Um \texttt{varchar} vai entre aspas na chamada.

Quando a consulta tem vários parâmetros (por exemplo, montar o padrão do
\texttt{like}), a professora considera o procedimento melhor que a
visão: ele recebe os valores e monta o \texttt{where} por dentro, em vez
de exigir um \texttt{where} de quem consulta a visão.

\subsection{Valor padrão}

\begin{lstlisting}
create procedure busca_livro2
@trecho varchar(60) = '%'
as
select codlivro, titulo from livro where titulo like @trecho
go

exec busca_livro2            -- sem argumento: traz todos
exec busca_livro2 'Dom%'
\end{lstlisting}

É por causa de valores padrão que alguns procedimentos do sistema
funcionam com ou sem parâmetro.

\subsection{Parâmetro de saída}

\begin{lstlisting}
create procedure titulo_do_livro
@codlivro int,
@titulo char(60) output       -- devolvido a quem chamou
as
select @titulo = titulo from livro where codlivro = @codlivro
go

declare @t char(60)
exec titulo_do_livro 10, @t output   -- output também na chamada
print @t
\end{lstlisting}

O parâmetro \texttt{output} é passado por referência. A ideia é não
mostrar o resultado na tela de dentro do procedimento e sim devolvê-lo
numa variável. A variável é atribuída no próprio \texttt{select}
(\texttt{select @titulo = titulo ...}). Quem chama declara uma variável
do mesmo tipo e escreve \texttt{output} também no \texttt{exec}. Sem o
\texttt{print}, nada aparece.

\subsection{Valor de retorno}

\begin{lstlisting}
create procedure conta_livros
@trecho varchar(60)
as
select codlivro, titulo from livro where titulo like '%' + @trecho + '%'
return @@rowcount
go

declare @r int
exec @r = conta_livros 'amor'
print @r
\end{lstlisting}

\begin{itemize}
	\item O \texttt{return} só devolve \textbf{inteiro}. Serve para status
	(sucesso ou falha), não para devolver o valor de uma coluna; para isso
	existe o \texttt{output}.
	\item O \texttt{return} encerra o procedimento ali mesmo, inclusive
	dentro de um \texttt{if}.
	\item O retorno é opcional: o procedimento funciona igual sem ele. Para
	usá-lo, declara-se uma variável e executa-se o bloco inteiro
	(\texttt{declare}, \texttt{exec @r = ...}, \texttt{print}).
	\item \texttt{print} aparece na aba \emph{Mensagens} do SSMS;
	\texttt{select @r} aparece na aba \emph{Resultados}, como uma
	consulta.
\end{itemize}

\subsection{Alterar e excluir}

\begin{lstlisting}
alter procedure busca_livro
@trecho varchar(60) = '%'
as
select codlivro, titulo from livro where titulo like '%' + @trecho + '%'
go

exec sp_helptext busca_livro           -- mostra o código
exec sp_rename 'busca_livro', 'procura_livro'
drop procedure procura_livro, busca_livro2
\end{lstlisting}

Como na visão, o \texttt{alter} substitui a definição e mantém as
permissões.

\section{Procedimentos com transação}

O uso principal na disciplina: encapsular uma transação num
procedimento que devolve 1 em caso de sucesso e 0 em caso de falha. O
padrão é o do \refch{cet615-dml}, com um \texttt{return} em cada ramo.

\begin{lstlisting}[caption={Cadastro de livro com exemplar: obra genérica e específica numa transação.}]
create procedure ins_livro_exemplar
@codlivro int,
@titulo char(60),
@isbn numeric(13,0),
@codautor int,
@codexemplar int
as
begin transaction
insert into livro (codlivro, titulo, isbn, codautor)
values (@codlivro, @titulo, @isbn, @codautor)
if @@rowcount > 0               /* livro inserido */
begin
   insert into exemplar (codexemplar, codlivro, estado)
   values (@codexemplar, @codlivro, 'N')   -- estado fixo: novo
   if @@rowcount > 0            /* exemplar inserido */
   begin
      commit transaction
      return 1
   end
   else
   begin
      rollback transaction
      return 0
   end
end
else
begin
   rollback transaction
   return 0
end
go

declare @ret int
exec @ret = ins_livro_exemplar 10, 'Dom Casmurro', 9788535910663, 1, 100
print @ret
\end{lstlisting}

\subsection{Como projetar os parâmetros}

A professora deu atenção especial a esse ponto:

\begin{itemize}
	\item \textbf{A chave vai uma vez só.} \texttt{@codlivro} serve para as
	duas tabelas. Passar o código duas vezes abre a chance de chegarem
	valores diferentes, um erro que ela provocou de propósito na transação
	sem procedimento.
	\item \textbf{O que é fixo não é parâmetro.} O tipo da especialização
	(0 para pessoa física, 0 para livro) e o estado inicial são escritos
	no próprio procedimento. Como parâmetro, só abririam espaço para erro.
	\item \textbf{Um procedimento por especialização.} Juntar o cadastro de
	pessoa física e jurídica num só (``a mania de querer aproveitar'')
	deixaria metade dos parâmetros nula em cada chamada. Para exclusão,
	um procedimento único faz sentido: basta o código, e o procedimento
	descobre o tipo.
	\item \textbf{A transação é a menor possível.} O que não precisa estar
	junto fica de fora. Por isso a avaliação diz que o vendedor já está
	cadastrado: cadastrar cliente é outra transação.
	\item \textbf{A ordem dos \texttt{insert} segue as chaves
	estrangeiras.} Primeiro a linha referenciada (a obra genérica), depois
	a que referencia (a específica).
\end{itemize}

\subsection{Por que um \texttt{return} em cada ramo}

Não dá para escrever um único \texttt{return @@rowcount} no fim, porque
o \texttt{@@rowcount} muda a cada comando. A documentação\sidecite{MSRowcount}
detalha: comandos DML gravam nele o número de linhas afetadas; comandos
como \texttt{print}, \texttt{begin transaction} e \texttt{commit
transaction} o zeram; atribuições simples o deixam em 1. Depois do
\texttt{commit}, ele já vale 0. A regra prática é testar o
\texttt{@@rowcount} logo depois do comando que interessa, ou guardá-lo
numa variável. Com \texttt{return} junto do \texttt{commit} ou do
\texttt{rollback}, o bloco \texttt{begin ... end} passa a ser
obrigatório em cada ramo.

\subsection{Procedimentos que chamam procedimentos}

Um procedimento pode executar outro. No exemplo da aula, um
procedimento de transação completa cadastra um cliente pessoa física e
duas contas (corrente e poupança) chamando os procedimentos menores e
testando o retorno de cada um:

\begin{lstlisting}
declare @retorno int
begin transaction
exec @retorno = ins_pessoa_fisica ...
if @retorno > 0
begin
   exec @retorno = ins_conta ... 0 ...   -- corrente
   if @retorno > 0
   begin
      exec @retorno = ins_conta ... 1 ...  -- poupança
      if @retorno > 0 commit transaction
      else rollback transaction
   end
   else rollback transaction
end
else rollback transaction
return @retorno
\end{lstlisting}

Aqui o \texttt{return @retorno} pode ficar no fim, porque
\texttt{@retorno} é uma variável comum e não muda sozinha. Modularizar
encurta o procedimento maior, e os menores já foram testados. Em
compensação, uma transação grande trava mais o banco, e a recomendação
continua sendo mantê-las enxutas.

\subsection*{Complemento: \texttt{rollback} dentro de procedimento aninhado}

A aula não tratou deste ponto, e ele afeta o exemplo acima no caminho de
falha. O SQL Server não tem transações aninhadas
independentes\sidecite{MSRollback}:

\begin{itemize}
	\item cada \texttt{begin transaction} soma 1 em
	\texttt{@@trancount}, e o \texttt{commit} de uma transação interna só
	subtrai 1. Quem efetiva de fato é o \texttt{commit} mais externo;
	\item um \texttt{rollback} sem ponto de salvamento desfaz tudo até o
	\texttt{begin transaction} mais externo e zera
	\texttt{@@trancount}, mesmo que tenha sido executado dentro do
	procedimento chamado;
	\item se o procedimento termina com \texttt{@@trancount} diferente do
	valor com que começou, o SQL Server emite uma mensagem informativa.
\end{itemize}

Na prática: se \texttt{ins\_conta} falha e executa \texttt{rollback}, a
transação inteira, inclusive o cliente já inserido, é desfeita ali. O
procedimento externo recebe 0 e executa o seu \texttt{rollback}, mas não
há mais transação aberta para desfazer. No caminho de sucesso, o exemplo
funciona como esperado. Uma forma de evitar o problema é o procedimento
interno só abrir transação quando \texttt{@@trancount} for 0.

\subsection*{Complemento: procedimentos no padrão SQL}

\textcite[seção 13.4]{Elmasri2011} apresenta os procedimentos
armazenados do padrão SQL/PSM (\emph{Persistent Stored Modules}), com
parâmetros \texttt{in}, \texttt{out} e \texttt{inout} e comandos de
controle de fluxo. Além do reuso por várias aplicações, que a aula
também cita, o livro aponta duas vantagens: menos tráfego entre cliente
e servidor, porque o processamento fica no servidor, e a possibilidade
de verificar regras complexas demais para as restrições declarativas. O
livro observa ainda que muitos SGBDs comerciais aceitam procedimentos
escritos em linguagens de uso geral. O \texttt{output} do T-SQL
corresponde ao \texttt{out} e ao \texttt{inout} do padrão.

\setchapterpreamble[u]{\margintoc}
\chapter{T-SQL: Gatilhos}
\labch{cet615-gatilhos}

Este capítulo cobre a apresentação sobre gatilhos da terceira
aula~\cite{BaiocoCET615}. Com eles, a disciplina fecha a sequência de
objetos do banco vista desde a primeira aula: tabelas, índices,
restrições, visões, procedimentos e gatilhos.

\section{O que é um gatilho}

Um \textbf{gatilho}\index{gatilho} (\emph{trigger}) é um conjunto de
comandos SQL ligado a uma tabela e executado automaticamente pelo SGBD
quando os dados dela mudam por \texttt{insert}, \texttt{update} ou
\texttt{delete}. Pode ser visto como um procedimento armazenado sem
parâmetros, que ninguém chama: ele dispara sozinho a cada comando a que
está associado. Gatilhos servem para verificar ou modificar dados, e não
para devolver dados ao usuário com \texttt{select}.

O paralelo com o procedimento armazenado ajuda: os dois são conjuntos
de comandos SQL, mas o procedimento não pertence a uma tabela e é
chamado com \texttt{exec}; o gatilho fica preso a uma tabela e a um
comando, não recebe parâmetros e não pode ser chamado.

Usos convencionais citados na aula:

\begin{itemize}
	\item manter integridade que os recursos declarativos não cobrem
	(\texttt{unique}, \texttt{primary key}, \texttt{not null},
	\texttt{foreign key});
	\item propagar modificações entre tabelas relacionadas;
	\item manter valores derivados, como o total de uma nota fiscal
	calculado a partir dos itens.
\end{itemize}

\subsection{Bancos de dados ativos}

Os gatilhos também implementam \textbf{bancos de dados ativos}: bancos
que, além de guardar dados e executar transações, reagem a eventos. As
reações são regras ativas no modelo \textbf{evento--condição--ação}:

\begin{lstlisting}[language={},numbers=none]
on evento    -- a modificação que dispara a regra
if condição  -- teste opcional
then ação    -- o que fazer
\end{lstlisting}

Na história contada em aula\sidecite{BaiocoCET615Aulas}, os gatilhos
existiam nos gerenciadores antes do padrão SQL e nasceram para checar
restrições de integridade, numa época em que declarar \texttt{primary
key} só documentava a coluna. Quando os fabricantes passaram a checar as
restrições sozinhos, o gatilho ganhou outro uso: tornar o banco ativo. O
usuário que escolhe uma opção do menu não sabe que outros comandos rodam
por trás; o projetista e o programador sabem.

\textcite[seção 26.1]{Elmasri2011} descreve o mesmo modelo: o evento
costuma ser uma atualização no banco, mas pode ser um evento de tempo
(``todo dia às 5h30'') ou externo; a condição é opcional; e a ação pode
ser uma sequência de comandos SQL, uma transação ou um programa externo.

\section{Funcionamento no SQL Server}

Cada comando SQL é uma transação. Quando a tabela tem um gatilho para
aquele comando, a transação é estendida com o código do gatilho. Um
\texttt{delete} passa a valer como:

\begin{lstlisting}[language={},numbers=none]
begin transaction
   grava o log
   remove as linhas
   executa o gatilho de exclusão
commit ou rollback
\end{lstlisting}

Por isso um \texttt{rollback} dentro do gatilho desfaz também o comando
que o disparou, e um \texttt{rollback} em quem executou o comando desfaz
também o que o gatilho fez.

Três consequências que a professora destacou:

\begin{itemize}
	\item O gatilho não tem \texttt{begin transaction}, porque continua a
	transação do comando, nem \texttt{commit}, porque pode haver uma
	transação externa com outros comandos ainda por executar.
	\item Ela sempre põe um ponto de \texttt{rollback} no gatilho. Com
	mais de uma operação, testa a primeira antes de fazer a segunda, para
	não gastar tempo com operações inúteis. Esse \texttt{rollback} desfaz
	também a transação externa inteira.
	\item Gatilho não alivia bloqueio. Trocar uma transação explícita por
	um gatilho não deixa nada mais leve: a transação continua do mesmo
	tamanho. O ganho é de consistência, porque as ações rodam sempre que
	alguém executa o comando, sem depender de lembrar de abrir uma
	transação.
\end{itemize}

\begin{figure}[h]
\centering
\begin{tikzpicture}[node distance=4mm]
	\node[caixa, fill=white] (cmd) {insert into exemplar ...};
	\node[caixaazul, below=of cmd, minimum width=50mm] (exec) {grava log e altera a tabela};
	\node[caixaazul, below=of exec, minimum width=50mm] (rest) {verifica restrições (PK, FK, check)};
	\node[caixa, below=of rest, minimum width=50mm, fill=bdroxo, text=white] (trg) {executa o gatilho};
	\node[caixaazul, right=6mm of trg, minimum width=18mm] (ins) {inserted\\deleted};
	\node[caixa, below=of trg, minimum width=50mm] (fim) {commit ou rollback};
	\draw[seta] (cmd) -- (exec);
	\draw[seta] (exec) -- (rest);
	\draw[seta] (rest) -- (trg);
	\draw[seta] (trg) -- (fim);
	\draw[seta, dashed] (ins) -- (trg);
	\node[draw=bdcinza, dashed, fit=(exec)(fim), inner sep=2.5mm,
		label={[rotulo]left:uma transação}] {};
\end{tikzpicture}
\caption{Um gatilho \texttt{after} no SQL Server roda dentro da mesma
transação do comando, depois das restrições, uma vez por comando.}
\end{figure}

\subsection{As tabelas \texttt{inserted} e \texttt{deleted}}

Durante o gatilho, o SQL Server oferece duas tabelas lógicas, com a
mesma estrutura da tabela do gatilho, que podem ser consultadas com
\texttt{select}:

\begin{description}
	\item[\texttt{inserted}:] as linhas novas de um \texttt{insert}, ou os
	valores novos de um \texttt{update}.
	\item[\texttt{deleted}:] as linhas removidas por um \texttt{delete}, ou
	os valores antigos de um \texttt{update}.
\end{description}

Um \texttt{update} funciona como uma exclusão seguida de uma inclusão: o
valor antigo vai para \texttt{deleted} e o novo, para \texttt{inserted}.
Não existe uma tabela \texttt{updated}. Comparar as duas mostra o que
mudou; a diferença \texttt{i.quantidade - d.quantidade} serve para os
dois sentidos, porque fica negativa quando a quantidade diminui. No
Oracle, os equivalentes são \texttt{new} e \texttt{old}.

A função \texttt{update(coluna)} indica se a coluna aparece no
\texttt{set} do comando. Ela não compara valores: devolve verdadeiro
mesmo que o valor não mude. Serve para um gatilho de alteração mais
pontual, que só age quando a coluna que interessa foi tocada, já que
disparar à toa custa tempo e bloqueio em outra tabela.

Se o mesmo valor de \texttt{inserted} é usado várias vezes nas ações,
a professora sugere carregá-lo uma vez numa variável em vez de repetir o
\texttt{select}.

\section{Criar, alterar e excluir}

\begin{lstlisting}
create trigger nome
on tabela
for {[insert] [,] [update] [,] [delete]}
as
comandos

alter trigger nome on tabela for ... as ...  -- substitui a definição
exec sp_rename 'nome', 'novo_nome'
drop trigger nome1, nome2
\end{lstlisting}

O gatilho aparece no SSMS dentro da tabela, na pasta \emph{Gatilhos}, e
\texttt{sp\_helptext} mostra o código, como em visões e procedimentos.
\texttt{drop table} apaga os gatilhos da tabela junto; visões e
procedimentos precisam de \texttt{drop} próprio.

\subsection{\texttt{for}, \texttt{after}, \texttt{before} e
\texttt{instead of}}

O slide só mostra \texttt{for}, a forma antiga. \texttt{after} veio com
a padronização, e no SQL Server os dois são iguais: primeiro executa o
comando, depois o gatilho. O padrão e o Oracle têm também
\texttt{before}, que inverte a ordem; o SQL Server não tem. Ele tem
\texttt{instead of}, que \emph{substitui} o comando. Muita gente usa
\texttt{instead of} para simular \texttt{before}, repetindo o comando
original dentro do gatilho; a professora chama isso de adaptação, não do
propósito do recurso.

\subsection{Projetar antes de escrever}

``Não adianta saber sintaxe e não saber projetar um gatilho''. Antes de
escrever, defina quais operações, em quais tabelas, o disparam, e quais
ações ele executa. Duas orientações da aula:

\begin{itemize}
	\item Um gatilho por operação. A sintaxe aceita \texttt{insert,
	update, delete} no mesmo gatilho, mas isso só é comum em gatilhos de
	verificação; para regras ativas, as ações mudam de uma operação para
	outra (o \texttt{delete} subtrai o que o \texttt{insert} somou), e um
	gatilho único com \texttt{if} para escolher o código não compensa.
	\item Pode haver mais de um gatilho para a mesma operação. Ações
	independentes podem ficar separadas; ações em que a ordem importa
	ficam no mesmo gatilho, porque não se controla a ordem de execução
	entre gatilhos.
\end{itemize}

\subsection{Gatilho ou transação explícita?}

Tudo o que um gatilho faz pode ser feito por um procedimento com
transação; o contrário nem sempre. A professora deu um teste que, nas
palavras dela, não está nos slides: pense no procedimento que faria a
transação. Se os parâmetros da primeira operação bastam para executar
as outras, as outras podem virar gatilho da primeira, porque esses
valores estarão em \texttt{inserted} ou \texttt{deleted}.

\begin{itemize}
	\item \textbf{Dá para ser gatilho:} inserir um item de nota recebe
	número da nota, produto e quantidade. Com isso, o gatilho busca o
	preço e atualiza o total da nota, sem parâmetro extra. O mesmo vale
	para baixar o estoque ao vender um item.
	\item \textbf{Não dá:} cadastrar um atendente, que é uma pessoa
	genérica mais um registro específico. A segunda inclusão precisa de
	salário e comissão, que a primeira não tem; só o código se
	reaproveita.
\end{itemize}

O método sugerido para o projeto é listar tudo o que exige mais de uma
operação e decidir, caso a caso, o que fica em transação explícita e o
que vai para gatilho.

\section{Exemplos}

O exemplo da aula usa produto, nota fiscal e item de nota, e mantém o
valor total da nota com três gatilhos. Aqui o mesmo raciocínio aparece
no banco de biblioteca: a tabela \texttt{livro} ganha uma coluna
derivada \texttt{qtd\_exemplares}, mantida pelos gatilhos da tabela
\texttt{exemplar}.

\begin{lstlisting}
alter table livro add qtd_exemplares int not null default 0
\end{lstlisting}

\begin{lstlisting}[caption={Inclusão: soma os exemplares novos de cada livro.}]
create trigger exemplar_inclusao
on exemplar for insert
as
update livro
set qtd_exemplares = qtd_exemplares + i.qtd
from livro l
     inner join (select codlivro, count(*) as qtd
                 from inserted
                 group by codlivro) i
     on l.codlivro = i.codlivro
\end{lstlisting}

\begin{lstlisting}[caption={Exclusão: desconta os exemplares removidos.}]
create trigger exemplar_exclusao
on exemplar for delete
as
update livro
set qtd_exemplares = qtd_exemplares - d.qtd
from livro l
     inner join (select codlivro, count(*) as qtd
                 from deleted
                 group by codlivro) d
     on l.codlivro = d.codlivro
\end{lstlisting}

\begin{lstlisting}[caption={Alteração: só age se o exemplar mudou de livro.}]
create trigger exemplar_alteracao
on exemplar for update
as
if update(codlivro)
begin
   update livro set qtd_exemplares = qtd_exemplares - d.qtd
   from livro l inner join (select codlivro, count(*) as qtd
                            from deleted group by codlivro) d
        on l.codlivro = d.codlivro

   update livro set qtd_exemplares = qtd_exemplares + i.qtd
   from livro l inner join (select codlivro, count(*) as qtd
                            from inserted group by codlivro) i
        on l.codlivro = i.codlivro
end
\end{lstlisting}

O \texttt{update ... from} com junção é a forma do T-SQL de alterar uma
tabela usando dados de outra. O agrupamento por \texttt{codlivro} faz o
gatilho funcionar mesmo quando um único comando mexe em vários
exemplares; o porquê está no complemento a seguir.

Os gatilhos da aula seguem outro padrão: atualizam a nota com
\texttt{where numeronota = (select numeronota from inserted)} e, se
nenhuma linha foi afetada (\texttt{@@rowcount = 0}), executam
\texttt{rollback}. Na resolução dos exercícios, a professora observou
que ``nenhuma linha afetada'' quer dizer ``não encontrou'', e não
encontrar nem sempre é erro: quem decide abortar é a lógica do
gatilho.

O exercício de fixação traz mais um uso: quando uma fatura é paga
(\texttt{if update(dtpagamento)}), o gatilho a apaga de \texttt{fatura}
e a insere numa tabela \texttt{faturapaga} com \texttt{insert ...
select * from inserted}. O mesmo recurso alimenta tabelas históricas ou
desnormalizadas para BI, juntando \texttt{inserted} com tabelas reais.

\subsection{Gatilhos e chaves estrangeiras}

Os testes da aula mostram uma interação importante. Inserir um item
para uma nota que não existe falha pela chave estrangeira, antes de o
gatilho rodar. E o gatilho que apaga os itens de uma nota excluída nunca
chega a rodar: a chave estrangeira do item impede o \texttt{delete} da
nota antes. A solução dos slides é remover essa chave estrangeira e
passar a verificar, no gatilho de inclusão, se a nota existe
(\texttt{if not exists (...)}, com \texttt{print} e
\texttt{rollback}).

A professora fez isso de propósito, para mostrar como a cascata era
feita antes de existir \texttt{on delete cascade}, e avisou que não é
para repetir. Tirou só a chave estrangeira do item para a nota, mantendo
a do produto, e pegou o nome gerado da restrição na própria mensagem de
erro (conferido com \texttt{sp\_helpconstraint}). Sem a chave, um item
para a nota 0, que não existe, passou: uma inconsistência que ela chamou
de inaceitável e que o novo gatilho de inclusão passou a barrar. No fim,
o \texttt{delete} da nota apagou também os itens. Cascata e integridade
funcionaram, só que por dois gatilhos, e o gerenciador hoje entrega as
duas coisas prontas.

\section*{Complementos da documentação}

\subsection*{O gatilho dispara uma vez por comando}

Os slides dizem que o gatilho se repete para cada linha afetada. Em
aula, a professora ressalvou que isso depende do fabricante, lembrou que
o SQL Server não tem o \texttt{for each row} do Oracle e deixou para a
turma verificar o caso de várias linhas num comando. No SQL Server não é
por linha: o gatilho DML dispara \emph{uma vez por comando},
mesmo que o comando afete várias linhas ou nenhuma, e as tabelas
\texttt{inserted} e \texttt{deleted} trazem todas as linhas
afetadas\sidecite{MSCreateTrigger}. A documentação recomenda escrever
gatilhos preparados para várias linhas\sidecite{MSTriggerInsertedDeleted}.

A consequência prática: \texttt{where numeronota = (select numeronota
from inserted)} só funciona quando o comando insere um item por vez. Um
\texttt{insert ... select} com itens de duas notas faz a subconsulta
devolver dois valores, e o comando falha. A forma segura é juntar com
\texttt{inserted} e agrupar, como nos exemplos acima.

\textcite[seção 26.1]{Elmasri2011} explica a origem da confusão: no
padrão SQL e no Oracle, o gatilho pode ser declarado \texttt{for each
row} (dispara por linha) ou não (dispara por comando). O SQL Server só
tem a segunda forma.

\subsection*{Quando as linhas já estão na tabela}

Os slides dizem que \texttt{inserted} contém as linhas que serão
inseridas, mas ainda não foram. No SQL Server, \texttt{for} e
\texttt{after} são sinônimos: o gatilho roda \emph{depois} que o comando
foi executado com sucesso, inclusive as verificações de restrições e as
ações em cascata\sidecite{MSCreateTrigger}. Quando o gatilho roda, as
linhas já estão na tabela, dentro da transação ainda aberta; é isso que
o \texttt{rollback} desfaz. É também por isso que a chave estrangeira
barra o comando antes do gatilho no teste da aula.

\subsection*{Alternativas para manter a chave estrangeira}

Remover a chave estrangeira para o gatilho de exclusão funcionar troca
uma restrição declarativa por código. Duas alternativas mantêm a
integridade referencial:

\begin{itemize}
	\item declarar a chave estrangeira com \texttt{on delete cascade}
	(\refch{cet615-ddl}): o próprio SGBD apaga os itens, sem gatilho;
	\item usar um gatilho \texttt{instead of delete} na nota, que roda
	\emph{no lugar} do comando, antes das restrições: ele apaga os itens
	e depois a nota. A documentação não permite esse tipo de gatilho numa
	tabela que tenha \texttt{on delete cascade}\sidecite{MSCreateTrigger}.
\end{itemize}

Os gatilhos \texttt{instead of} também são a forma de tornar atualizável
uma visão que não seria (\refch{cet615-visoes}).

\subsection*{\texttt{commit} dentro do gatilho}

Em aula, a professora lembrou, como recordação incerta, que o SQL Server
ignora um \texttt{commit} no gatilho e só avisa. O teste mostra outra
coisa. Num gatilho de inclusão com \texttt{commit}, o SQL Server 2022
confirmou a inclusão e o que o gatilho fez, e em seguida deu o erro 3609
(\emph{the transaction ended in the trigger; the batch has been
aborted}): o resto do lote não rodou, e nenhuma transação externa
poderia mais desfazer aquelas linhas. É mais um motivo para a regra da
aula: gatilho não tem \texttt{commit}.

\subsection*{Gatilhos encadeados e recursivos}

Um aluno perguntou se um \texttt{update} dentro de um gatilho pode
disparar outro gatilho. Pode, se as opções do servidor permitirem. A
professora separou dois casos: o encadeamento, que vem ligado e em si não
causa problema; e a recursão, quando a cadeia volta a disparar um
gatilho anterior, que vem bloqueada porque, sem condição de parada, vira
laço infinito.

Os padrões batem com o SQL Server 2022: a opção de servidor
\texttt{nested triggers} vem ligada, e a opção de banco
\texttt{recursive\_triggers}, desligada\sidecite{MSCreateTrigger}. O
detalhe é o alcance: \texttt{recursive\_triggers} desligada só impede a
recursão \emph{direta}, o gatilho de T1 alterando T1. A indireta (T1
dispara T2, que altera T1) só é barrada desligando \texttt{nested
triggers}, o que barra todo encadeamento. Com ela ligada, um laço
indireto para no limite de 32 níveis, com erro.

\subsection*{Outros pontos}

\begin{itemize}
	\item Gatilhos que devolvem resultados com \texttt{select} estão
	marcados para remoção em versão futura.
	\item \texttt{truncate table} não dispara gatilho de exclusão.
	\item Sobre a ordem entre gatilhos da mesma operação: a aula diz que
	não há controle. Há um pouco: \texttt{sp\_settriggerorder} define o
	primeiro e o último gatilho \texttt{after}; os do meio rodam em
	ordem indefinida. A recomendação da aula, juntar no mesmo gatilho o
	que depende de ordem, continua a mais segura.
	\item Além dos gatilhos de dados, o SQL Server tem gatilhos de DDL,
	disparados por \texttt{create}, \texttt{alter}, \texttt{drop},
	\texttt{grant} e outros, no banco (\texttt{on database}) ou no
	servidor (\texttt{on all server}), e gatilhos de logon. A aula não
	tratou deles\sidecite{MSCreateTrigger}.
	\item Um gatilho mantém os bloqueios da transação enquanto roda, então
	deve ser curto. A documentação sugere sair logo quando o comando não
	afetou linha nenhuma: \texttt{if (rowcount\_big() = 0) return}.
\end{itemize}

\setchapterpreamble[u]{\margintoc}
\chapter{T-SQL: Controle de Dados (DCL)}
\labch{cet615-dcl}

Este capítulo cobre a última apresentação da
disciplina~\cite{BaiocoCET615}: controle de acesso, cópia de
segurança e agendamento de tarefas. A aula tratou o tema como uma visão
geral: quem já trabalha como DBA conhece, e quem não conhece fica sabendo
o que existe para pesquisar depois. Não houve demonstração ao vivo; as
telas do SSMS são as dos slides.

\section{Por que usar o controle do SGBD}

Muitos sistemas fazem o controle de acesso só na aplicação: uma tabela
de usuários e senhas dentro do banco e uma conexão que pode tudo. A
professora chamou isso de ``bem temeroso'' por causa da \textbf{injeção
de SQL}\index{injeção de SQL} (\emph{SQL injection}): com trechos de
comando digitados no campo de login, forçando erros não tratados,
alguém descobre um acesso válido, e por trás está uma conexão sem
restrição. A recomendação é usar o controle de acesso do próprio SGBD,
que nos fabricantes estabelecidos é bem elaborado e serve a aplicações em
qualquer linguagem.

\section{Três níveis de acesso}

Para ler um dado, um usuário precisa passar por três portas:

\begin{enumerate}
	\item \textbf{Servidor:} um \emph{login} permite conectar à instância.
	\item \textbf{Banco:} um \emph{user} (usuário do banco), ligado ao
	login, permite entrar em um banco específico.
	\item \textbf{Objetos:} \emph{permissões} dizem o que o usuário pode
	fazer com cada tabela, visão ou procedimento.
\end{enumerate}

Um login sem usuário no banco conecta ao servidor mas não abre o banco;
um usuário sem permissão na tabela abre o banco mas não lê a tabela. O
\texttt{sa} (\emph{system administrator}) é um login com todas as
permissões; um login comum pode servir só para conectar.

\begin{figure}[h]
\centering
\begin{tikzpicture}
	\node[caixa, fill=white, minimum width=20mm] (p) at (0,0) {pessoa ou\\aplicação};
	\node[draw=bdroxo, fill=bdroxoclaro, rounded corners, minimum width=74mm, minimum height=36mm] (srv) at (5.6,0) {};
	\node[rotulo, anchor=north west] at (srv.north west) {servidor: login};
	\node[draw=bdazul, fill=bdazulclaro, rounded corners, minimum width=54mm, minimum height=24mm] (bd) at (6.3,-0.4) {};
	\node[rotulo, anchor=north west] at (bd.north west) {banco: usuário};
	\node[caixa, fill=white, minimum width=36mm] (obj) at (6.6,-0.8) {objetos: permissões\\(select, execute ...)};
	\draw[seta] (p) -- node[rotulo, above] {conecta} (srv.west);
\end{tikzpicture}
\caption{Os três níveis de acesso, de fora para dentro.}
\end{figure}

\section{Logins e usuários}

\begin{lstlisting}
use master
go
create login leitor_app
with password = 'Senha#Forte2026', default_database = biblioteca
go
alter login leitor_app with password = 'Outra#Senha2026'
drop login leitor_app

use biblioteca
go
create user app_biblioteca for login leitor_app
drop user app_biblioteca
\end{lstlisting}

O \texttt{create login} roda no \texttt{master}, onde fica o catálogo
de logins, e o \texttt{create user}, no banco de destino. O usuário não
tem senha: a senha é do login, e o usuário só o mapeia para dentro do
banco. Sem \texttt{default\_database}, o banco padrão do login é o
\texttt{master}, o que a professora considera ruim. Os procedimentos antigos (\texttt{sp\_addlogin},
\texttt{sp\_adduser}, \texttt{sp\_droplogin} e outros) foram substituídos
por esses comandos a partir do SQL Server 2005. Tudo isso também pode ser
feito pelo SSMS: a pasta \emph{Security} do servidor tem logins e papéis
de servidor, e a pasta \emph{Security} de cada banco tem usuários,
papéis e esquemas.

\section{Papéis (\emph{roles})}

Controle de acesso é projeto, ``não é assim do nada''. Um
\textbf{papel}\index{role} é um grupo de usuários com outro nome: agrupa
quem precisa das mesmas permissões.
Quem entra no papel herda as permissões dele, o que simplifica a
administração. Há dois tipos:

\begin{description}
	\item[Papéis de servidor:] agrupam logins e dão permissões
	administrativas. O \texttt{sysadmin} pode tudo na instância; outros
	cobrem partes do trabalho, como \texttt{securityadmin} (logins),
	\texttt{serveradmin} (opções do servidor) e \texttt{dbcreator}
	(criação de bancos). A professora sugere que o DBA os use para delegar
	tarefas, ou para separar as próprias, em vez de fazer tudo como
	\texttt{sa}. Desde o SQL Server 2012 é possível criar papéis de
	servidor próprios (\texttt{create server role}); ela prefere combinar
	os papéis prontos.
	\item[Papéis de banco:] agrupam usuários de um banco. Há papéis
	prontos, como \texttt{db\_owner} (controle total do banco),
	\texttt{db\_datareader}, \texttt{db\_datawriter} e
	\texttt{db\_ddladmin}, que permite criar, alterar e excluir objetos e
	não deve ir para o usuário final. É possível criar novos com
	\texttt{create role}. Todo usuário pertence automaticamente ao papel
	\texttt{public}.
\end{description}

\begin{lstlisting}
alter server role sysadmin add member leitor_app   -- login
alter server role sysadmin drop member leitor_app

use biblioteca
create role atendentes
alter role atendentes add member app_biblioteca    -- usuário
alter role atendentes drop member app_biblioteca
drop role atendentes
\end{lstlisting}

\subsection*{Complemento: o que o papel \texttt{public} permite}

O slide diz que o \texttt{public} tem permissões limitadas a leitura de
dados, e em aula a professora alertou que, por isso, um usuário recém
criado já poderia fazer \texttt{select}. A documentação diz que o
\texttt{public} não tem permissões próprias sobre os dados: ele lê
metadados do sistema\sidecite{MSDatabaseRoles}, e o usuário herda dele só o que alguém conceder ao
\texttt{public}. No teste, um usuário novo, só no \texttt{public},
recebeu \emph{The SELECT permission was denied} ao consultar uma tabela.
O alerta continua útil: o que for concedido ao \texttt{public} vale para
todos.

\section{Permissões}

Permissões definem quem pode fazer o quê, em três níveis:

\begin{description}
	\item[Servidor:] vêm principalmente dos papéis de servidor; desde o
	SQL Server 2012 também podem ser dadas diretamente a um login.
	\item[Banco:] criar objetos, fazer backup do banco ou do log. Vêm dos
	papéis de banco ou de \texttt{grant} direto.
	\item[Objeto:] consultar, inserir, alterar e excluir dados de uma
	tabela, executar um procedimento.
\end{description}

\begin{lstlisting}
-- banco
grant create table, create view to app_biblioteca
revoke create table, create view from app_biblioteca

-- objeto
grant select, insert on emprestimo to atendentes
grant execute on ins_livro_exemplar to atendentes
revoke insert on emprestimo from atendentes
deny delete on emprestimo to atendentes
\end{lstlisting}

Na prática, as permissões de banco são de DDL e as de objeto, de DML.
As de objeto chegam até a coluna, mas fazer isso com \texttt{grant} um a
um é trabalhoso, e é aí que os papéis ajudam.

\begin{description}
	\item[\texttt{grant}] concede a permissão.
	\item[\texttt{revoke}] retira uma permissão concedida \emph{ou
	negada} antes; o usuário volta ao estado neutro e pode ainda receber
	a permissão por um papel.
	\item[\texttt{deny}] nega explicitamente. A negação tem precedência:
	quem tem \texttt{deny} não recebe a permissão nem por papel.
\end{description}

A professora deu dois usos para o \texttt{deny}: negar temporariamente
sem retirar a permissão, durante uma investigação, por exemplo; e criar
um papel só de negações para os dados sensíveis, no qual entra todo
usuário novo. Combinar um papel que concede com outro que nega também
evita criar um papel para cada combinação, como um papel ``só
\texttt{insert}''.

O slide escreve \texttt{revoke ... to}; a documentação aceita
\texttt{to} e \texttt{from} nesse comando\sidecite{MSRevoke}, e
\texttt{from} lê melhor.

Este é o mecanismo por trás da ``segurança por funcionalidade'' dos
capítulos de visões e procedimentos: em vez de dar \texttt{select} nas
tabelas, dá-se \texttt{select} na visão ou \texttt{execute} no
procedimento, e o usuário nunca toca as tabelas. A professora insistiu
nisso: uma visão que já restringe linhas e colunas, ou um procedimento
que faz uma transação em várias tabelas, dá ao usuário exatamente a
funcionalidade de que ele precisa. Ela acha esses recursos pouco usados,
enquanto se tenta resolver tudo na interface, com brechas.

\subsection*{Complemento: \texttt{grant all}}

Em aula, \texttt{grant all} no banco foi comparado a pôr o usuário no
\texttt{db\_owner}. A documentação diz outra coisa: \texttt{all} está
obsoleto, existe só por compatibilidade e não concede todas as
permissões. No banco, equivale a \texttt{backup database},
\texttt{backup log}, \texttt{create database}, \texttt{create default},
\texttt{create function}, \texttt{create procedure}, \texttt{create
rule}, \texttt{create table} e \texttt{create view}\sidecite{MSGrant}.
O \texttt{db\_owner} tem o controle total do banco.

\subsection*{Complemento: exceções à precedência do \texttt{deny}}

A documentação registra uma exceção mantida por compatibilidade: um
\texttt{deny} na tabela inteira não prevalece sobre um \texttt{grant} em
uma coluna específica dela\sidecite{MSRevoke}. Também registra que a
forma \texttt{revoke all} está obsoleta; o certo é revogar permissões
específicas.

\subsection*{Complemento: propagação de permissões}

\textcite[seção 24.2]{Elmasri2011} trata o controle de acesso
discricionário do padrão SQL. Quem recebe uma permissão \texttt{with
grant option} pode repassá-la a outros. Ao revogar, o SGBD precisa
retirar também as permissões que foram repassadas a partir dela. No SQL
Server, isso exige \texttt{revoke grant option for ... cascade}; sem o
\texttt{cascade}, o \texttt{revoke} falha quando a permissão foi
concedida com \texttt{grant option}\sidecite{MSRevoke}. O mesmo capítulo
apresenta o controle de acesso baseado em papéis (seção 24.3), que é o
modelo dos \emph{roles} do SQL Server.

\section{Esquemas}

Um \textbf{esquema}\index{esquema (schema)} (\emph{schema}) agrupa
objetos de um banco que têm o mesmo dono, um usuário ou um papel. Os
donos passam a ser donos de esquemas, e não de cada objeto. Permissões
podem ser dadas no esquema inteiro. Dentro de um esquema, nomes de
objetos não se repetem. O nome completo de um objeto é
\texttt{servidor.banco.esquema.objeto}. Trocar o dono de um esquema não
muda esse nome.

O uso que a professora destacou: quando se decide dar acesso direto às
tabelas, sem visão nem procedimento, concede-se uma vez no esquema que
agrupa as tabelas, em vez de repetir o \texttt{grant} em cada uma. Ela
vê o melhor uso num banco central de empresa com departamentos bem
separados.

\begin{lstlisting}
grant select on schema :: dbo to atendentes
\end{lstlisting}

\subsection*{Complemento: comandos de esquema}

Em aula, ficou a dúvida se há comandos para criar esquemas e mover
objetos; há. O teste abaixo rodou no SQL Server 2022 e terminou com
\texttt{dbo.livro}, \texttt{acervo.livro} e \texttt{acervo.exemplar}
no mesmo banco:

\begin{lstlisting}
-- create schema precisa ser o primeiro comando do lote
create schema acervo
go
-- mesmo nome de dbo.livro, em outro esquema
create table acervo.livro (codlivro int)
-- move a tabela para o esquema acervo
alter schema acervo transfer dbo.exemplar
-- troca o dono do esquema
alter authorization on schema :: acervo to app_biblioteca
\end{lstlisting}

O teste confirma a regra do slide: o nome só precisa ser único dentro
do esquema, e dois esquemas do mesmo banco podem ter tabelas com o mesmo
nome. O esquema pertence a um banco e não agrupa objetos de bancos
diferentes.

\section{Cópia de segurança e recuperação}

Backup não é copiar os arquivos \texttt{.mdf}, \texttt{.ldf} e
\texttt{.ndf} pelo sistema operacional. Ele é feito de dentro do
servidor, porque a estrutura está no catálogo, que também precisa de
cópia. Como na modelagem, a estratégia depende dos requisitos e do valor
da informação: o banco precisa ficar no ar
24 horas por dia? Quanto tempo de parada se aceita depois de uma falha?
As possibilidades citadas na aula:

\begin{description}
	\item[Espelhamento (\emph{mirroring}):] os dados ficam duplicados em
	dois servidores, de preferência em lugares diferentes; se um falha, o
	outro assume em segundos ou minutos. Para a professora, é ``o melhor
	dos mundos'' e a única saída para bancos que mudam o tempo todo, como
	os de um banco comercial; uma loja pode viver com instantâneos ou
	backup completo.
	\item[Instantâneos (\emph{snapshots}):] cópias do banco em intervalos
	curtos.
	\item[Backup e restore:] cópias para restauração posterior.
\end{description}

\subsection{Tipos de backup}

\begin{description}
	\item[Completo (\emph{full}):] o banco inteiro. É o mais lento.
	\item[Diferencial:] só as páginas alteradas desde o último completo.
	\item[Do log de transações:] só o log.
\end{description}

Duas estratégias básicas:

\begin{itemize}
	\item \textbf{Só o banco}, pelo menos diariamente. O que mudou desde o
	último backup se perde numa falha.
	\item \textbf{Banco e log.} O completo é menos frequente (por exemplo,
	semanal), e o log é copiado várias vezes por dia. Na recuperação,
	restaura-se o último completo e depois todos os backups de log, em
	ordem cronológica. O diferencial encurta essa sequência.
\end{itemize}

O exemplo da aula para o backup só do banco: um comércio que funciona
oito horas por dia pode aceitar ficar o dia descoberto e fazer o backup
completo no fim do expediente. O objetivo de qualquer estratégia é
cobrir o maior período possível de alterações.

O \texttt{master} deve ter backup regular e também logo depois de criar,
alterar ou excluir bancos, logins, usuários, papéis ou permissões.
``Regular'' quer dizer que, sem essas mudanças, ele fica parado; depois
delas, o backup é imediato, ou se perde o catálogo. A professora
recomenda também um backup completo do banco de usuário depois de
mudanças de estrutura.

\begin{lstlisting}
backup database biblioteca to disk = 'C:\backup\biblioteca.bak'
backup log biblioteca to disk = 'C:\backup\biblioteca_log.trn'

restore database biblioteca from disk = 'C:\backup\biblioteca.bak'
   with norecovery          -- ainda virão logs
restore log biblioteca from disk = 'C:\backup\biblioteca_log.trn'
   with recovery            -- último passo: o banco volta ao ar
\end{lstlisting}

No SSMS: botão direito no banco $\rightarrow$ \emph{Tasks} $\rightarrow$
\emph{Back Up} ou \emph{Restore}. A aula cita duas novidades do SQL Server
2005: a restauração parcial, que deixa o banco disponível assim que o
grupo de arquivos primário é restaurado, e o backup do log em paralelo
com o backup completo.

\subsection*{Complemento: backup com o banco no ar}

Em aula, foi dito que backup e restauração podem exigir tirar o banco do
ar. No SQL Server, os backups completo, diferencial e de log rodam com o
banco on-line e em uso; o custo é de leitura em disco e desempenho. No
teste, um \texttt{backup database} rodou com uma conexão usando o
banco. Quem costuma deixar o banco indisponível é a restauração, e é por
isso que a restauração parcial foi tratada como novidade.

Também apareceu a ideia de que o \emph{checkpoint} é o momento de
esvaziar o log. São coisas diferentes: o checkpoint grava em disco as
páginas alteradas que estão na memória. No modelo de recuperação
completo, o espaço do log só é liberado depois do backup do log; no
simples, depois do checkpoint\sidecite{MSTransactionLog}.

\subsection*{Complemento: modelos de recuperação}

O que pode ser copiado depende do \textbf{modelo de recuperação} do
banco\sidecite{MSRecoveryModels}:

\begin{description}
	\item[Simples:] não permite backup de log. O SQL Server reaproveita o
	espaço do log sozinho. Numa falha, perde-se o que mudou desde o
	último backup.
	\item[Completo (\emph{full}):] exige backup de log, e o log cresce até
	que ele seja feito. Permite restaurar até um ponto exato no tempo,
	por exemplo antes de um erro de um usuário.
	\item[\emph{Bulk-logged}:] variação do completo que registra pouco as
	cargas em massa; não permite restaurar até um ponto no tempo.
\end{description}

As edições Enterprise e Standard usam o modelo completo por padrão; a
Express usa o simples. A estratégia ``banco e log'' da aula só funciona
no modelo completo ou no \emph{bulk-logged}. Isso explica também o
alerta do \refch{cet615-sqlserver} sobre logs que crescem demais: no
modelo completo, sem backup de log, o log não é truncado.

\subsection*{Complemento: o espelhamento está obsoleto}

A documentação marca o espelhamento de banco para remoção em versão
futura e recomenda os grupos de disponibilidade \emph{Always On} para
alta disponibilidade\sidecite{MSMirroring}. O espelhamento também só
funciona com o modelo de recuperação completo. Já o instantâneo de banco
do SQL Server é uma visão somente leitura, na mesma instância e
dependente do banco de origem: não substitui o backup.

\section{Agendamento de tarefas (\emph{jobs})}

Um \textbf{job}\index{job} é uma tarefa administrativa, com um ou mais
passos, executada automaticamente em datas e horários definidos: backups,
manutenção de índices e outras rotinas do DBA. Também pode ser criado
sem agendamento, para execução manual. Os jobs dependem do serviço
\emph{SQL Server Agent} e ficam guardados no banco \texttt{msdb}
(\refch{cet615-sqlserver}). No SSMS: \emph{SQL Server Agent}
$\rightarrow$ \emph{Jobs} $\rightarrow$ \emph{New Job}, com as abas
\emph{Steps} (passos) e \emph{Schedules} (agendamento). Os passos podem
usar os próprios comandos \texttt{backup} e \texttt{restore}. Mesmo sem
agendamento, o job pronto facilita a vida do DBA.

Com os jobs, a professora encerrou o conteúdo da disciplina. Segurança e
administração devem voltar com mais profundidade em outras disciplinas
do curso.

\setchapterpreamble[u]{\margintoc}
\chapter{Avaliações}
\labch{cet615-avaliacoes}

A disciplina teve duas avaliações individuais, entregues pelo Moodle como
arquivo TXT ou PDF; a nota final é a média das duas, e a média para
aprovação é 7. A professora pediu respostas assertivas: só os comandos
pedidos, sem dados de teste e sem alterar o esquema físico fornecido.

\section{Primeira avaliação}

Entregue até 14/11/2026, às 9h. As duas questões usam o esquema de um
sistema de venda de obras literárias usadas, o mesmo do mapeamento do
\refch{cet615-conceitos}:

\begin{itemize}
	\item \texttt{cliente(login, senha, nome, endereco,
	qtd\_obras\_vendidas, qtd\_obras\_compradas)};
	\item \texttt{obras\_literarias(codobra, nome, editora, preco, tipo,
	login\_vendedor, data)}, com \texttt{tipo} 0 para livro e 1 para
	periódico, e \texttt{login\_vendedor} referenciando \texttt{cliente};
	\item \texttt{livro(codobra, ISBN)} e \texttt{periodico(codobra,
	numero, volume)}, cada um com \texttt{codobra} referenciando a obra;
	\item \texttt{compra(codcompra, data, login\_comprador, codobra)}, com
	um índice único em \texttt{codobra}: cada obra é comprada no máximo
	uma vez.
\end{itemize}

\begin{figure*}[h]
\centering
\begin{tikzpicture}[
	tabela/.style={rectangle split, rectangle split parts=2, draw=bdroxo, thick,
		rectangle split part fill={bdroxo, white}, font=\scriptsize, align=left,
		inner sep=2pt, text width=31mm},
	fk/.style={-{Stealth[length=2mm]}, draw=bdazul, thick}]
	\node[tabela] (cli) at (0,0) {\textcolor{white}{\textbf{cliente}}
		\nodepart{two} \underline{login}\\ senha, nome, endereco\\ qtd\_obras\_vendidas\\ qtd\_obras\_compradas};
	\node[tabela] (obr) at (5.2,0) {\textcolor{white}{\textbf{obras\_literarias}}
		\nodepart{two} \underline{codobra}\\ nome, editora, preco\\ tipo\\ login\_vendedor (FK)\\ data};
	\node[tabela] (com) at (10.4,0) {\textcolor{white}{\textbf{compra}}
		\nodepart{two} \underline{codcompra}\\ data\\ login\_comprador (FK)\\ codobra (FK, única)};
	\node[tabela] (liv) at (3.0,-3.2) {\textcolor{white}{\textbf{livro}}
		\nodepart{two} \underline{codobra} (FK)\\ ISBN};
	\node[tabela] (per) at (7.4,-3.2) {\textcolor{white}{\textbf{periodico}}
		\nodepart{two} \underline{codobra} (FK)\\ numero, volume};
	\draw[fk] (obr.west) -- node[rotulo, above] {vende} (cli.east);
	\draw[fk] (com.west) -- node[rotulo, above] {obra} (obr.east);
	\draw[fk] (com.north) -- ++(0,0.4) -| node[rotulo, above, pos=0.25] {comprador} (cli.north);
	\draw[fk] (liv.north) -- (liv.north |- obr.south);
	\draw[fk] (per.north) -- (per.north |- obr.south);
\end{tikzpicture}
\caption{Tabelas do esquema da primeira avaliação, desenhadas a partir do
esquema lógico. As setas vão da chave estrangeira para a tabela
referenciada; chaves primárias sublinhadas.}
\end{figure*}

\subsection{Questão 1 (4,0): visão}

Pedido: uma visão que, para cada periódico já comprado, mostre o login e
o nome do vendedor, o código e o nome da obra, a data em que foi posta à
venda e a data da compra; e um exemplo de consulta à visão.

\begin{lstlisting}
create view periodicos_comprados
as
select c.login login_vendedor, c.nome nome_vendedor,
       o.codobra, o.nome nome_obra,
       o.data data_venda, cp.data data_compra
from cliente c inner join obras_literarias o
         on c.login = o.login_vendedor
         inner join periodico p
         on o.codobra = p.codobra
         inner join compra cp
         on o.codobra = cp.codobra
go

select * from periodicos_comprados
where login_vendedor = 'ana'
order by data_compra
\end{lstlisting}

Decisões:

\begin{itemize}
	\item A junção com \texttt{periodico} restringe a obra aos periódicos;
	\texttt{o.tipo = 1} daria o mesmo resultado com dados consistentes. A
	junção com \texttt{compra} restringe às obras já compradas.
	\item As duas colunas \texttt{nome} e as duas \texttt{data} recebem
	apelidos, porque teriam o mesmo nome na visão.
	\item A visão traz só o que foi pedido; o filtro e a ordenação ficam na
	consulta de exemplo, como recomendado no \refch{cet615-visoes}.
\end{itemize}

\subsection{Questão 2 (6,0): procedimento armazenado}

Pedido: um procedimento para a transação de cadastro de livros, que
retorne 1 em caso de sucesso e 0 caso contrário, com um exemplo de
execução que mostre o retorno. Os vendedores já estão cadastrados.

\begin{lstlisting}
create procedure ins_livro
@codobra int,
@nome char(40),
@editora char(40),
@preco numeric(6,2),
@login_vendedor char(8),
@data date,
@isbn numeric(13,0)
as
begin transaction
insert into obras_literarias
values (@codobra, @nome, @editora, @preco, 0, @login_vendedor, @data)
if @@rowcount > 0
begin
      insert into livro
      values (@codobra, @isbn)
      if @@rowcount > 0
      begin
            commit transaction
            return 1
      end
      else
      begin
            rollback transaction
            return 0
      end
end
else
begin
      rollback transaction
      return 0
end
go

declare @ret int
exec @ret = ins_livro 10, 'Dom Casmurro', 'Garnier', 25.00, 'ana',
            '2026-11-07', 9788535910663
print @ret
\end{lstlisting}

Decisões, seguindo o \refch{cet615-procedimentos}:

\begin{itemize}
	\item A transação tem as duas operações do cadastro: a obra genérica e
	o livro específico, nessa ordem por causa da chave estrangeira.
	\item O código da obra é passado uma vez só, e o \texttt{tipo} é fixo
	(0), não parâmetro.
	\item Os tipos dos parâmetros são os mesmos das colunas do esquema
	físico.
	\item A contagem \texttt{qtd\_obras\_vendidas} do vendedor ficou de
	fora. Na apresentação da avaliação, a professora pediu a transação
	mais compacta possível e anunciou que a segunda avaliação usaria o
	mesmo esquema. Totais derivados como esse são o exemplo clássico de
	gatilho (\refch{cet615-gatilhos}).
\end{itemize}

\subsection{Verificação}

As duas respostas foram testadas no SQL Server 2022 com dados de exemplo
antes da entrega. A visão devolveu só os periódicos comprados, deixando
de fora um periódico não comprado e um livro comprado. O procedimento
retornou 1 num cadastro válido. Também retornou 0 com código de obra
repetido e com vendedor inexistente, sem deixar nada gravado e sem
transação aberta. Nos casos de falha, o SQL Server mostra a mensagem do
erro (violação de chave primária ou estrangeira) antes do retorno 0. É o
comportamento esperado do padrão com \texttt{@@rowcount}.

\section{Segunda avaliação}

\textit{Será incluída depois da entrega, prevista para 21/11/2026, às 9h.}


%----------------------------------------------------------------------------------------

\backmatter % Denotes the end of the main document content
\setchapterstyle{plain} % Output plain chapters from this point onwards

%----------------------------------------------------------------------------------------
%	BIBLIOGRAPHY
%----------------------------------------------------------------------------------------

% The bibliography needs to be compiled with biber using your LaTeX editor, or on the command line with 'biber main' from the template directory

\printbibliography[heading=bibintoc, title=Referências] % Add the bibliography heading to the ToC and output the bibliography

%----------------------------------------------------------------------------------------
%	GLOSSARY
%----------------------------------------------------------------------------------------

% The glossary needs to be compiled on the command line with 'makeglossaries main' from the template directory

\setglossarystyle{listgroup} % Set the style of the glossary (see https://en.wikibooks.org/wiki/LaTeX/Glossary for a reference)
\printglossary[title=Glossário, toctitle=Glossário] % Output the glossary, 'title' is the chapter heading for the glossary, toctitle is the table of contents heading

%----------------------------------------------------------------------------------------
%	INDEX
%----------------------------------------------------------------------------------------

% The index needs to be compiled on the command line with 'makeindex main' from the template directory

\printindex % Output the index

%----------------------------------------------------------------------------------------
%	BACK COVER
%----------------------------------------------------------------------------------------

% If you have a PDF/image file that you want to use as a back cover, uncomment the following lines

%\clearpage
%\thispagestyle{empty}
%\null%
%\clearpage
%\includepdf{cover-back.pdf}

%----------------------------------------------------------------------------------------

\end{document}
